\section{Roadmap to solved long-phrase transcription}
\label{sec:roadmap}

This section is the working plan for making Curtis able to produce long, trustworthy audio--video--transcription--score evidence from Alan's actual violin practice. It is intentionally part of the paper rather than a separate planning note. Whenever Curtis makes real progress toward full-session active-practice measurement, accurate transcription, score matching, heat maps, repertoire updates, or Curtis-level observation, this PDF should be updated with the new state, the acceptance evidence, and the remaining blocker.

\subsection{Current tracker state}

The tracker was refreshed through the 2026-05-18 source-lane eighth-note extension, the owner-media A4-audit pass, the notation rendering passes, the 04:43 PM EDT phrase-continuity gate, the 05:19 PM EDT visible-mismatch quarantine, the 06:14 PM EDT source-crop reverification queue, the 06:32 PM EDT packaged packet fallback, the 07:56 PM EDT source/audio acceptance attempt, the 2026-05-19 six-note audio gate, the 2026-05-19 failed-note raw-audio probe, Alan's 2026-05-19 live rejection of the re-promoted May 3 measure-16 display, the 2026-05-20 review-queue repeat-suppression update, and the 2026-05-20 score-transcription lane gate. The system has a real media ledger, a deployed active-practice scan route, a working owner-sync path for captured media, a visible 100-point transcription-completion tracker with precise weighted progress, stored user-correction benchmarks, a persistent truth workbench for accepted/rejected/pending audio-score-transcription items, a local source-backed Wieniawski score PDF, a nine-note source-backed Wieniawski symbolic score opening phrase, source-target triage queues, an unaccepted score-glyph verification queue, source review packets, hidden YIN fast-transition candidates, a phrase-audit packet generator, adjacent mining, source-audio rescanning, temporal-continuity checks, a raw continuity-probe path, a failed-adjacent-note raw probe, a source-crop reverification gate that can accept a rejected lane only when full-context source review and paired-audio review agree, a gold-review queue that hides exact normalized MIDI patterns, rejected clip/score fingerprints, overlapping rejected clip windows, and already accepted source areas, plus separate visible \texttt{transcription-alan} and \texttt{score-transcription} review lanes. \texttt{score-transcription} packets test exact visual notation copying--pitch, octave, accidentals, key signature, durations, rests, beams, tuplets, stems, noteheads, spacing, slurs/ties, and source range--without counting as \texttt{transcription-alan} practice evidence. The Wieniawski opening source map is \texttt{A5 G5 F5 A5 G5 F5 A5 G\#5 F5}; the previous \texttt{D6 C6 Bb5 D6 C6 Bb5 D6 C\#6 B5} and lower-octave \texttt{D5 C5 Bb4 D5 C5 Bb4 D5} maps are rejected and covered by regression tests. The 2026-05-03 Scherzo-Tarantelle five-note phrase \texttt{A\#4 D5 C5 A\#4 D5}, displayed as \texttt{Bb4 D5 C5 Bb4 D5}, remains rejected. The old A4/A5 single-note anchors remain rejected. The May 3 measure-16 display \texttt{Eb5 Eb5 C5 Eb5 Eb5} is now rejected too: Alan confirmed the audio is measure 16, but the transcription should begin \texttt{Eb D C}, and the displayed source crop is not the correct measure-16 snippet. Curtis now has no accepted visible score match, no accepted longer source-lane expansion, and a stricter evidence gate preventing internal staff indices, filename-based measure labels, source crops alone, or repeated rejected review candidates from becoming accepted matches. The notation renderer still uses local SMuFL/ABC engraving safeguards for accepted snippets, but future phrase growth is limited to ranges whose source crop, measure, note sequence, box centers, paired audio, transcription, continuity, and truth review all agree.

\begin{table}[htbp]
\centering
\tiny
\setlength{\tabcolsep}{3pt}
\renewcommand{\arraystretch}{0.88}
\caption{Progress tracker baseline for the long-phrase transcription goal.}
\label{tab:roadmap_baseline}
\begin{tabularx}{\textwidth}{>{\bfseries}p{0.22\textwidth}X}
\toprule
Tracker item & Baseline state \\
\midrule
Goal & Produce compact daily evidence groups: original score snippet, generated transcription, local audio, local video, heat-map context, and specific observation. \\
Current accepted transcription & No May 3 score-linked transcription is currently accepted. The re-promoted five-note measure-16 display is revoked because Alan's live review found that the audio is measure 16, the transcription should begin \texttt{Eb D C} rather than \texttt{Eb Eb C}, and the displayed source crop is not the correct measure-16 snippet. Audio-checked D/A fragments remain useful note evidence, but they are not score-location matches. \\
Current accepted score match & No May 3 score match is currently accepted. The old tight crop, red-bracket crop, six-/seven-/eight-note dependent ranges, the 2026-05-03 \texttt{Bb4 D5 C5 Bb4 D5} crop, the old A4/A5 single-note anchors, the previous D--B-flat--G--D score extraction, and the stitched A4 continuation remain rejected or blocked fixtures. The current blocker is creating one new source crop whose boxed score notes, measure location, generated transcription, and paired audio agree before any compact score--clip--audio--transcription group can return. The 951 local score glyph candidates, 214 staff-position notehead hypotheses, 8 staff review packets, and source-verification targets are implementation queues, not accepted score evidence. \\
Current archive coverage & 276 h 23 min uploaded practice video across 52 strict-ledger videos and 45 practice days; 2 h 24 min checked; 2 h 21 min measured active violin; 270 h 37 min provisional archive estimate; full chronological coverage still pending. \\
Current estimate & 223 h 5 min total active practice estimated from checked windows; provisional and sample-biased until the full chronological scan covers the archive. \\
Current user burden to remove & Alan is still forced to inspect whether a note is real, whether a score box points to the correct pitch, whether audio matches notation, and whether practice time means playing rather than uploaded duration. The review queue should no longer make him accept or reject the same exact normalized MIDI pattern, clip window, score image, or shifted overlapping clip area repeatedly. \\
First tracker update & This PDF now defines the full implementation plan and becomes the required progress record for future Curtis work. \\
\bottomrule
\end{tabularx}
\end{table}

\subsubsection*{Progress log}

\paragraph{2026-05-13 11:57 AM EDT code update} Backend coverage ledger, evidence-correction route, and benchmark-fixture scaffold added. Checked video time now includes sampled non-playing windows; total practice time remains measured active playing only. Live deployed coverage now reports 213 h 23 min uploaded video, 2 h checked video, 53 min 14 s measured active practice, and 211 h 23 min unchecked. Wrong score-note corrections can be stored as rejected benchmark cases. Full archive scan and long-phrase transcription remain unsolved.
\paragraph{2026-05-13 12:33 PM EDT code update} Low-cost active-practice scan route implemented. The new scan persists active violin-playing sub-intervals from stored violin-positive media samples, keeps checked non-playing samples separate from practice time, exposes \url{/api/curtis/active-practice-scan}, adds an active-scan run endpoint, queues remaining strict-ledger windows for owner media capture, and feeds persisted active intervals into daily records and the coverage ledger. Local unit tests verify active sub-interval extraction, negative-sample withholding, and coverage accounting.
\paragraph{2026-05-13 12:59 PM EDT deployed scan update} The active-scan endpoint was deployed and run against the live Curtis state. It persisted 214 active sub-intervals from 80 stored media-sample results, representing 41 min 42 s of scan-derived active audio, and queued 250 strict-ledger windows for later owner-media capture. The coverage ledger reports 1 h 57 min 24 s measured active practice from 2 h checked and a 208 h 46 min provisional archive estimate. A regression fix now caps active practice by checked source coverage and makes active-scan results supersede noisy transcription-duration evidence where scan results exist. The full 213 h 23 min archive scan and long-phrase transcription remain unsolved.
\paragraph{2026-05-13 03:26 PM EDT queue update} The owner-media sync now consumes the active-scan pending-window queue before choosing expansion or blind sample windows. The active-scan status endpoint can expose a bounded pending-window preview, and the run endpoint returns a compact scan/coverage payload rather than the full operations document. Regression tests now verify that queued active-practice windows are selected first and that pending-window previews preserve total pending counts while limiting returned rows. This does not solve long-phrase transcription; it closes the next M1 plumbing gap so the chronological archive scan can keep advancing without Alan manually choosing windows.
\paragraph{2026-05-13 03:37 PM EDT live owner-sync update} A one-window owner-sync pass uploaded one cached violin-positive 5/3 sample, \texttt{Njh8\_zq9\_DM-10725} from \texttt{*10725-10815}. The follow-up active scan selected that new sample and recorded one new active interval. Live coverage is now 81 active-scan sample results, 215 active sub-intervals, 42 min 9 s scan-derived active audio, 2 h 1 min 30 s checked video, 1 h 58 min 54 s measured active practice, and a 208 h 49 min provisional estimate. The first chronological pending queue item remains \texttt{violin 1} at \texttt{*0-90}; this run advanced cached-media coverage, not the chronological queue.
\paragraph{2026-05-13 04:05 PM EDT completion tracker and queue-only update} Curtis now exposes a bottom-page transcription completion tracker backed by a 100-point implementation gate checklist. The tracker reports archive coverage, accepted long phrases, exact score windows, scan queue state, completed gates, remaining gates, and the next concrete action. A queue-only owner-sync mode was added so the active-scan pending-window list can be consumed without falling back to cached positives, expansion windows, or blind sampling. A queue-only run uploaded \texttt{otHfHMgDo2g-0} from \texttt{violin 1} at \texttt{*0-90}; the sample was classified as \texttt{not\_violin\_or\_unclear} and therefore did not add practice time. The follow-up active scan reports 82 sample results, 72 active-violin samples, 10 checked non-violin samples, 215 active sub-intervals, 2 h 1 min 30 s checked video, 1 h 58 min 54 s measured active practice, and a 208 h 49 min provisional estimate. The next chronological pending window is \texttt{violin 1} at \texttt{*90-180}.
\paragraph{2026-05-13 05:44 PM EDT queue resume and practice-time ceiling update} The next queue-only owner-sync attempt captured \texttt{violin 1} windows \texttt{*90-540}. The wrapper timed out after the uploads, revealing that cached pending-window media was excluded from queue-only resume. The owner-sync helper now resumes cached queue windows, including cached non-playing windows, and tests cover that path. The follow-up active scan reports 87 sample results, 77 active-violin samples, 10 checked non-violin samples, 238 active sub-intervals, 2 h 3 min checked video, 1 h 55 min 30 s measured active practice, and a 200 h 23 min provisional estimate. Practice time is now hard-capped by checked source coverage, so unchecked transcription remnants cannot make active practice exceed checked video. The next chronological pending window is \texttt{violin 1} at \texttt{*540-630}.
\paragraph{2026-05-13 06:01 PM EDT chronological queue update} A one-window queue-only owner-sync pass uploaded \texttt{otHfHMgDo2g-540} from \texttt{violin 1} at \texttt{*540-630}. The sample was classified as \texttt{violin\_positive} with a sampler score of 68.6. The follow-up active scan selected that sample, recorded 4 new active sub-intervals, and reports 88 sample results, 78 active-violin samples, 10 checked non-violin samples, 242 active sub-intervals, 2 h 4 min 30 s checked video, 1 h 57 min measured active practice, and a 200 h 32 min provisional estimate. The live accounting check passes because active practice remains below checked video. The next chronological pending window is \texttt{violin 1} at \texttt{*630-720}.
\paragraph{2026-05-13 06:26 PM EDT chronological queue update} A one-window queue-only owner-sync pass uploaded \texttt{otHfHMgDo2g-630} from \texttt{violin 1} at \texttt{*630-720}. The sample was classified as \texttt{violin\_positive} with a sampler score of 60.3. The follow-up active scan selected that sample, recorded 5 new active sub-intervals, and reports 89 sample results, 79 active-violin samples, 10 checked non-violin samples, 247 active sub-intervals, 2 h 6 min checked video, 1 h 58 min 30 s measured active practice, and a 200 h 41 min provisional estimate. The live accounting check passes because active practice remains below checked video. The next chronological pending window is \texttt{violin 1} at \texttt{*720-810}.
\paragraph{2026-05-13 07:55 PM EDT chronological queue update} A one-window queue-only owner-sync pass uploaded \texttt{otHfHMgDo2g-720} from \texttt{violin 1} at \texttt{*720-810}. The sample was classified as \texttt{violin\_positive} with a sampler score of 55.2. The follow-up active scan selected that sample, recorded 3 new active sub-intervals, and reports 90 sample results, 80 active-violin samples, 10 checked non-violin samples, 250 active sub-intervals, 2 h 6 min checked video, 2 h measured active practice, and a 203 h 14 min provisional estimate. The live accounting check passes because active practice remains below checked video. The next chronological pending window is \texttt{violin 1} at \texttt{*810-900}.
\paragraph{2026-05-13 09:25 PM EDT scan and benchmark update} Queue-only owner sync consumed the next two chronological \texttt{violin 1} windows, \texttt{*810-900} and \texttt{*900-990}. Both were classified as \texttt{violin\_positive}. The follow-up active scan reports 92 sample results, 82 active-violin samples, 10 checked non-violin samples, 257 active sub-intervals, 2 h 10 min checked video, 2 h 3 min measured active practice, and a 201 h 7 min provisional archive estimate. The next chronological pending window is \texttt{violin 1} at \texttt{*990-1080}. The same pass stored three real user-reviewed score-match failures as benchmark corrections: \texttt{A4} displayed against a B-position score crop, \texttt{A4} displayed against a \texttt{G4} crop, and \texttt{A5} displayed against a \texttt{G5} crop. These benchmark fixtures moved the implementation tracker from 29 percent to 31 percent. Exact score windows and accepted long phrases remain at 0.
\paragraph{2026-05-13 10:05 PM EDT scan and exact-percent update} Queue-only owner sync consumed \texttt{violin 1} windows \texttt{*990-1080} and \texttt{*1080-1170}. Both were classified as \texttt{violin\_positive}. The follow-up active scan reports 94 sample results, 84 active-violin samples, 10 checked non-violin samples, 269 active sub-intervals, 2 h 13 min checked video, 2 h 6 min measured active practice, and a 201 h 24 min provisional archive estimate. The next chronological pending window is \texttt{violin 1} at \texttt{*1170-1260}. The rounded implementation label remains 31 percent, but exact weighted progress moved to 31.13 percent and is now displayed so small coverage gains are visible. Exact score windows and accepted long phrases remain at 0.
\paragraph{2026-05-13 11:14 PM EDT scan and precision update} Queue-only owner sync consumed \texttt{violin 1} windows \texttt{*1170-1260} and \texttt{*1260-1350}. The first was classified as \texttt{violin\_positive}; the second was classified as \texttt{not\_violin\_or\_unclear}, which adds checked non-playing coverage without adding practice time. The follow-up active scan reports 96 sample results, 85 active-violin samples, 11 checked non-violin samples, 283 active sub-intervals, 2 h 16 min checked video, 2 h 7 min measured active practice, and a 199 h 19 min provisional archive estimate. The next chronological pending window is \texttt{violin 1} at \texttt{*1350-1440}. The implementation tracker now computes exact weighted progress from unrounded gate points, so the live label should read 31.128 percent after deployment. Exact score windows and accepted long phrases remain at 0.
\paragraph{2026-05-13 11:38 PM EDT scan update} Queue-only owner sync consumed \texttt{violin 1} windows \texttt{*1350-1440} and \texttt{*1440-1530}. Both were classified as \texttt{violin\_positive}. The follow-up active scan selected both samples, recorded 11 new active sub-intervals, and reports 98 sample results, 87 active-violin samples, 11 checked non-violin samples, 294 active sub-intervals, 2 h 19 min checked video, 2 h 10 min measured active practice, and a 199 h 37 min provisional archive estimate. The next chronological pending window is \texttt{violin 1} at \texttt{*1530-1620}. Exact weighted implementation progress is now 31.131 percent. Exact score windows and accepted long phrases remain at 0.
\paragraph{2026-05-13 11:55 PM EDT scan update} Queue-only owner sync consumed \texttt{violin 1} windows \texttt{*1530-1620} and \texttt{*1620-1710}. The first was classified as \texttt{not\_violin\_or\_unclear}; the second was classified as \texttt{violin\_positive}. The follow-up active scan selected both samples, recorded 6 new active sub-intervals, and reports 100 sample results, 88 active-violin samples, 12 checked non-violin samples, 300 active sub-intervals, 2 h 22 min checked video, 2 h 12 min measured active practice, and a 197 h 40 min provisional archive estimate. The next chronological pending window is \texttt{violin 1} at \texttt{*1710-1800}. Exact weighted implementation progress is now 31.134 percent. Exact score windows and accepted long phrases remain at 0.
\paragraph{2026-05-14 12:09 AM EDT scan update} Queue-only owner sync consumed \texttt{violin 1} windows \texttt{*1710-1800} and \texttt{*1800-1890}. Both were classified as \texttt{violin\_positive}. The follow-up active scan selected both samples, recorded 13 new active sub-intervals, and reports 102 sample results, 90 active-violin samples, 12 checked non-violin samples, 313 active sub-intervals, 2 h 25 min checked video, 2 h 15 min measured active practice, and a 197 h 59 min provisional archive estimate. The next chronological pending window is \texttt{violin 1} at \texttt{*1890-1980}. Exact weighted implementation progress is now 31.136 percent. Exact score windows and accepted long phrases remain at 0.
\paragraph{2026-05-14 12:35 AM EDT scan update} Queue-only owner sync consumed \texttt{violin 1} windows \texttt{*1890-1980} and \texttt{*1980-2070}. Both were classified as \texttt{violin\_positive}. The follow-up active scan selected both samples, recorded 9 new active sub-intervals, and reports 104 sample results, 92 active-violin samples, 12 checked non-violin samples, 322 active sub-intervals, 2 h 28 min checked video, 2 h 18 min measured active practice, and a 198 h 18 min provisional archive estimate. The next chronological pending window is \texttt{violin 1} at \texttt{*2070-2160}. Exact weighted implementation progress is now 31.139 percent. Exact score windows and accepted long phrases remain at 0.
\paragraph{2026-05-14 03:00 AM EDT scan update} Queue-only owner sync consumed \texttt{violin 1} windows \texttt{*2070-2160} and \texttt{*2160-2250}. Both were classified as \texttt{violin\_positive}. The follow-up active scan folded both samples into the checked ledger and reports 106 sample results, 94 active-violin samples, 12 checked non-violin samples, 331 active sub-intervals, 2 h 31 min checked video, 2 h 21 min measured active practice, and a 198 h 36 min provisional archive estimate. The next chronological pending window is \texttt{violin 1} at \texttt{*2250-2340}. Exact weighted implementation progress is now 31.142 percent. Exact score windows and accepted long phrases remain at 0.
\paragraph{2026-05-14 03:23 AM EDT one-measure accelerator update} Curtis no longer hardcodes accepted long phrases to zero. The tracker now derives accepted long-phrase count from daily match groups that pass all of the following gates: score location verified, exact score-location status, at least a five-note phrase run, local playable media, and not merely a single-pitch anchor. Regression tests now prove that a verified one-measure symbolic-score phrase with media increments the long-phrase count, while single-note anchors, unverified reference matches, and matches without media do not. This is implementation progress toward a faster path, not solved musical evidence yet: the live percent remains 31.142 percent until Curtis has a verified symbolic score measure or more scanned archive coverage. The fastest route to visible musical progress is now to verify one Scherzo-Tarantelle symbolic score measure and run the existing phrase matcher across the hidden detected note series already stored for 2026-05-03.
\paragraph{2026-05-14 04:32 AM EDT source-score acceleration update} Curtis now carries the Wieniawski Scherzo-Tarantelle solo score PDF as a local source asset rather than depending on weak scanned crops or a network fetch during matching. The score-asset layer reports whether a local source PDF is present, the daily score-reference audit counts local source PDFs separately from symbolic notes, and the 100-point tracker awards source-truth progress for a local source score without counting it as an accepted phrase. Regression tests verify that this raises the score-source readiness state while keeping accepted long phrases, exact score windows, and symbolic score notes at 0. This is the necessary acceleration layer for long-phrase work: the next step is no longer finding a score image, but converting one local source-score measure into verified symbolic notes and running the existing phrase matcher over stored detected note series.
\paragraph{2026-05-14 05:23 AM EDT live symbolic-measure update} The deployed Curtis API reported 43.642 percent exact progress, 1 provisional measure match, 1 provisional score-aligned window, 0 accepted long phrases, 1 local score source, and 8 symbolic score notes counted across the two Wieniawski practice-day targets. This state is now superseded: the later source-truth correction rejected the provisional \texttt{D A\# G D} source-score claim because the local IMSLP solo part did not support that score-note sequence. The useful implementation gain from this pass was the source-symbolic note gate and score/audio/video/transcription lane, not the specific provisional score phrase.
\paragraph{2026-05-14 05:40 AM EDT source-target accounting update} The deployed Curtis API reported 45.142 percent exact progress after the tracker began counting configured source score targets from the reference-training state instead of only local PDF files. That accounting layer remains useful, but the musical evidence count from this pre-correction state is superseded by the 07:52 AM source-truth correction. This makes the progress meter more honest and removes a hidden undercount, but it does not solve phrase transcription; the next implementation step is extending corrected accepted symbolic measure gates into longer matched note runs.
\paragraph{2026-05-14 06:52 AM EDT phrase-candidate ranking update} The deployed Curtis API reported 47.642 percent exact progress after the matcher stopped returning the first scan-order matches and instead ranked all detected-series matches by exact score status, score-derived status, distinct pitch-class content, and run length. That state exposed 12 pitch-sequence groups, 1 phrase candidate, 1 provisional measure match, 1 provisional score-aligned window, and 0 accepted long phrases. The highest-ranked pending phrase candidate was \texttt{D D D\# D D A\# G}. This candidate was not accepted score evidence; it was a focused source-score verification target. The later strict evidence gate supersedes this provisional count.

\paragraph{2026-05-14 07:08 AM EDT tracker visibility update} Curtis now returns the highest-ranked pending phrase target as structured tracker state and displays it in the implementation card detail as \texttt{pending: D D D\# D D A\# G}. The visible interface therefore shows not only that one phrase candidate exists, but also which detected pitch-class sequence should be source-score verified next. This does not add accepted musical evidence or change the 47.642 percent full-solve score; it removes a product friction point by making the next verification target visible without presenting it as a solved score match.

\paragraph{2026-05-14 07:52 AM EDT source-truth correction update} The Wieniawski symbolic source opening was corrected against the local IMSLP solo-violin score. The old accepted \texttt{D A\# G D} match was demoted because the source phrase is \texttt{D C Bb D C Bb}; it can remain only as an unverified reference-sequence candidate, not accepted score evidence. The deployed provisional groups matched detected pitch classes \texttt{A\# D C A\#} to source score pitch classes \texttt{A\# D C A\#}, displayed a source-score crop for that phrase, and retained generated notation only as fallback. The tracker then reported 48.642 percent, 2 provisional measure matches, 2 provisional score-aligned windows, 12 pitch-sequence groups, 1 pending phrase candidate, 3 score targets, 1 local score PDF, 6 unique symbolic score notes, 1 source snippet candidate, and 0 accepted long phrases. The later strict evidence gate supersedes these provisional score counts.

\paragraph{2026-05-14 08:29 AM EDT first source-backed phrase update} The Scherzo-Tarantelle source excerpt was extended by one verified score note, from \texttt{D C Bb D C Bb} to \texttt{D C Bb D C Bb D}, using a widened crop from the local IMSLP solo-violin PDF rather than the annotated teaching image. That seven-note source phrase let the stored 2026-05-03 detected series match the contiguous five-note run \texttt{A\# D C A\# D} against source score notes \texttt{Bb4 D5 C5 Bb4 D5}. The deployed API then reported 55.642 percent exact progress, 1 provisional long-phrase candidate, 2 provisional measure matches, 2 provisional score-aligned windows, 12 pitch-sequence groups, 7 unique symbolic score notes, and 1 source snippet candidate. The useful gain was the source-backed path; the later visual review rejected the crop as accepted score evidence, so the strict current state is again 0 accepted long phrases.

\paragraph{2026-05-14 01:50 PM EDT score-coordinate heat-map update} Provisional symbolic score phrase matches were connected to score-coordinate heat-map fragments. The 2026-05-03 Scherzo-Tarantelle phrase \texttt{A\# D C A\# D} was mapped to \texttt{mm. 2--4} with source score notes \texttt{Bb4 D5 C5 Bb4 D5}, local video, local audio, and an IMSLP score crop. The deployed API reported 58.136 percent exact progress, 1 provisional long-phrase candidate, 2 provisional measure matches, 2 provisional score-aligned windows, 2 provisional score-coordinate heat-map fragments, 12 pitch-sequence groups, 7 unique symbolic score notes, and 1 source snippet candidate. This was a heat-map and evidence-organization advance; the later strict evidence gate blocks those fragments from counting as accepted evidence.

\paragraph{2026-05-14 02:15 PM EDT source-verification target update} Curtis now separates the next longer note-run target from accepted score evidence. The tracker exposes one source-verification target, \texttt{D D D\# D D A\# G}, only because it is a seven-note run from local playable practice media with enough pitch diversity to be worth checking against the score. It remains reference-audio evidence, not accepted score evidence. The source target appears in the tracker and API so the next pass can verify exact local IMSLP score notes and coordinates instead of slowly searching from scratch or accidentally promoting a wrong crop. The deployed tracker state then included provisional score artifacts, but the current strict evidence gate counts 0 accepted long phrases, 0 accepted measure matches, 0 exact score-aligned windows, 0 score-coordinate heat-map fragments, and 0 verified source snippets until source review passes.

\paragraph{2026-05-14 02:51 PM EDT source-target triage update} Curtis now checks four-note-or-longer source targets against the verified symbolic score map before promotion. The current top target \texttt{D C G A} was compared against the verified source excerpt \texttt{D C A\# D C A\# D}. The best contiguous overlap is two notes, so the target is recorded as \texttt{source\_score\_sequence\_not\_found} and remains reference-audio evidence only. The tracker then reported 58.633 percent exact progress, 3 source-target score checks, 0 verified source targets, and the next action: extend the verified IMSLP symbolic score map before promoting this sequence. This improves the long-phrase pipeline by making measure-sized candidates fail fast and explicitly instead of appearing as misleading score matches.

\paragraph{2026-05-14 03:49 PM EDT score-map workbench update} Curtis now has an explicit source-map acceleration layer for the local IMSLP Wieniawski solo part. A deterministic extraction tool renders page 2 of the local score PDF, detects staff-line groups, removes staff lines, extracts connected score glyph candidates, and writes a JSON verification queue. The current queue contains 779 unaccepted score glyph candidates across 7 staves. These candidates are not pitch labels, not accepted score notes, and not displayable score evidence. They exist to make the next work cycle faster: verify candidate glyphs against the local PDF, convert only verified notes into MusicXML, then run source-target matching against verified symbolic notes. The tracker now reports 59.633 percent exact progress after awarding score-truth credit for the verification queue while preserving 0 verified source targets for the current \texttt{D C G A} source target.

\paragraph{2026-05-14 04:24 PM EDT note-hypothesis acceleration update} The source-map workbench now filters the local IMSLP glyph queue into likely notehead candidates and assigns staff-position pitch hypotheses using the G-minor key signature. The current queue contains 172 unaccepted score-note hypotheses across 7 staves in addition to the 779 raw glyph candidates. These hypotheses are explicitly not accepted score evidence: they use staff position only, may include false positives, and cannot create a score snippet, pitch label, or match card until source-reviewed into MusicXML. The implementation value is speed. The next pass can review a note-labeled queue instead of starting from anonymous glyph boxes, which should make extending the Scherzo-Tarantelle symbolic map to full measures and longer phrases much faster. The tracker now reports 61.133 percent exact progress after awarding limited score-truth credit for the note-hypothesis queue while preserving 0 verified source targets for the current \texttt{D C G A} source target.

\paragraph{2026-05-14 05:31 PM EDT staff-review packet update} The score-map detector was changed from a single-threshold staff finder to a multi-threshold merged staff finder after review showed that the previous queue skipped a source staff. Page-2 coverage now reports 951 unaccepted score glyph candidates and 214 unaccepted notehead hypotheses across 8 staves. The same pass writes 8 staff-level review-packet images under the local score assets folder, with each packet showing the scanned IMSLP staff plus numbered hypothesis boxes. The packets are not accepted score evidence and do not change any match card by themselves; they exist to make the next MusicXML pass faster and less error-prone by turning source review into a bounded staff-by-staff checklist. The deployed tracker now reports 64.135 percent exact progress after awarding score-truth credit for complete source-review packet generation while preserving 0 verified source targets for the current \texttt{D D D\# D D A\# G} source target, whose verified symbolic-score overlap is still only 1/7.

\paragraph{2026-05-14 06:15 PM EDT score/transcription visual-agreement update} The public 2026-05-03 match had a real failure: the score panel showed a broad opening source crop even though the claimed transcription phrase was the narrower five-note slice \texttt{Bb4 D5 C5 Bb4 D5}. Curtis then added a separate score-display gate: a score panel could render an actual source crop only when the crop's reference start and end exactly matched the phrase, and broad covering crops were withheld. The symbolic matcher also began requiring MIDI-level equality between detected notes and parsed score notes, not only pitch-class similarity, before it could create a score-visual lock. The deployed tracker reported 64.878 percent exact progress and provisional score locks, but later visual review showed the crop still was not trustworthy enough; the current strict gate counts this phrase as pending, not accepted.

\paragraph{2026-05-14 07:04 PM EDT actual-source score-gate update} The previous visual-agreement fix still left one unacceptable loophole: a generated notation SVG could appear in the accepted score panel even though Alan asked for the actual piece. Curtis therefore required a score panel to use a non-generated source-score crop for the same reference range. The five-note 2026-05-03 Scherzo-Tarantelle match gained an exact-range source crop for \texttt{Bb4 D5 C5 Bb4 D5}; generated notation remained only as the transcription rendering. That deployment expected provisional actual-source score progress, but the later musician review rejected the crop and the current strict evidence gate blocks it from creating an accepted measure, long-phrase count, exact score-aligned window, score-coordinate heat-map fragment, actual-source score lock, or verified source snippet.

\paragraph{2026-05-14 09:45 PM EDT score-spelling visual-context update} The latest reported mismatch exposed another gate that had been too weak: a transcription could be score-spelled as \texttt{Bb4 D5 C5 Bb4 D5} in data while the rendered notation panel omitted the source key signature, making the first and fourth notes look like unqualified natural-position B notes to a reader. Curtis now carries the source key signature into score-matched transcription rendering and requires both \texttt{scoreVisualRangeAgreement} and \texttt{scoreSpellingAgreement} before an actual-source score crop can count as accepted score evidence. The Scherzo-Tarantelle match remains a five-note fragment, but the visible transcription and visible score must now agree in musical spelling context rather than only MIDI or pitch class. This is a stricter acceptance gate, not a new long-phrase solve.

\paragraph{2026-05-15 12:29 PM EDT exact-note truth-gate update} The reported failure mode was not only visual range: a score match could still be treated as acceptable if the pitch class agreed while the written note/register did not. The truth workbench now stores MIDI sequences for detected, accepted, and score notes; accepted score-truth items reject audio-score mismatches at exact note-and-octave level; and score-truth items cannot be accepted with score notes that omit octave/register. Single-note score anchors also require an explicit exact-note verification flag before the UI can render them. Regression tests now cover accepted A4--D5 truth, A4-versus-A5 rejection, score truth without octave rejection, exact pitch-anchor display, and visual-only pitch-anchor rejection. This does not create a new accepted phrase; it prevents Curtis from ever calling a pitch-class-only or visually guessed crop solved evidence.

\paragraph{2026-05-14 10:52 PM EDT visual-crop demotion update} Alan's follow-up review showed that the notes visible in the 2026-05-03 source crop still did not match the transcription. Curtis now rejects the Scherzo-Tarantelle source crops as review material unless the visible noteheads, register, written spelling, and exact range are independently verified against the transcription. The five-note 5/3 phrase remains a symbolic phrase candidate, but it no longer creates an accepted score panel, accepted measure match, accepted long phrase, exact score-aligned window, score-coordinate heat-map fragment, score-visual lock, actual-source score lock, or verified source snippet. This is a regression lock: wrong score evidence is removed rather than cosmetically relabeled.

\paragraph{2026-05-15 12:45 AM EDT truth-workbench update} Curtis now has a persistent truth-workbench API and compact implementation-status surface for accepted, pending, and rejected audio-score-transcription evidence. A truth item can store the source clip, practice day, piece, detected notes, accepted notes, score notes, score source, score location, score image, and gate state. Accepted truth items are rejected at write time if the accepted audio note sequence and score note sequence disagree. The backend and frontend score gates also now require explicit visible-note-sequence verification before a score panel can count as a score-visual lock, actual-source score lock, accepted measure match, or accepted long phrase. This prevents the exact failure Alan found: a visually wrong score crop cannot count just because the metadata says the pitch sequence should match. It still has 0 accepted long phrases, 0 accepted measure matches, 0 exact score-aligned windows, 0 score-coordinate heat-map fragments, 0 score-visual locks, 0 actual-source score locks, and 0 accepted truth clips.

\paragraph{2026-05-15 03:20 PM EDT exact-source phrase attempt} Curtis temporarily promoted one actual-source Scherzo-Tarantelle phrase crop, \texttt{Bb4 D5 C5 Bb4 D5} from the local IMSLP page-2 solo-violin scan, aligned to the symbolic score range \texttt{mm. 2--4}. The attempt boxed five matched noteheads and stored source page pixel coordinates and exact visible note-plus-octave sequence. This pass was useful because it proved the crop extraction path could render actual score material, but it was later rejected by visual review because the score notes and boxes were not trustworthy enough to count as accepted evidence.

\paragraph{2026-05-15 03:40 PM EDT superseded deployment verification} The deployed Curtis site served cache version \texttt{20260515-exact-source-phrase} on deployment \texttt{4fb78ce1}. That deployment still showed one score group for 2026-05-03, with detected notes \texttt{A\#4 D5 C5 A\#4 D5}, score-spelled transcription \texttt{Bb4 D5 C5 Bb4 D5}, source score notes \texttt{Bb4 D5 C5 Bb4 D5}, the exact-source crop asset, local video sample \texttt{Njh8\_zq9\_DM-9465}, and local audio from 6.618--7.674 s inside that sample. This state is now superseded by the 05:00 PM EDT strict evidence-gate pass because the displayed score crop did not survive musician visual review. The live practice-time surface reported 234 h 57 min uploaded archive video, 2 h 28 min checked video, 2 h 21 min measured active violin, 223 h 5 min estimated total active practice, and 232 h 28 min unchecked video. The estimate remains provisional until the full chronological active-practice scan covers every archive window.

\paragraph{2026-05-15 05:00 PM EDT strict evidence-gate update} Curtis now rejects the 2026-05-03 Scherzo-Tarantelle source crop as accepted evidence unless the full truth chain passes: source note boxes centered on the visible noteheads, visible score notes matching the claimed written notes and octaves, transcription notes matching the paired audio, transcription notes matching the score notes, and an accepted truth-workbench item tying the audio, transcription, and score together. Backend, frontend, tracker, and tests now require those fields before a score snippet, accepted measure match, exact score-aligned window, actual-source score lock, or long-phrase count can appear. The deployed tracker reports 43.126 percent exact weighted progress, 0 accepted long phrases, 0 accepted measure matches, 0 exact score-aligned windows, 0 actual-source score locks, and 0 verified source snippets until a crop survives this stricter review.

\paragraph{2026-05-15 04:01 PM EDT practice-time API unification} The daily-record API now consumes the same active-practice coverage ledger used by \texttt{/api/curtis/active-practice-coverage} before publishing top-level and per-day practice totals. This removes the stale split where one public route could show older transcription-derived totals while the coverage ledger showed 2 h 28 min checked video, 2 h 21 min measured active violin, 223 h 5 min estimated total active practice, and 232 h 28 min unchecked video. A regression test now verifies that daily records inherit canonical checked, active, estimate, unmeasured, and per-day scan totals from the coverage ledger.

\paragraph{2026-05-15 06:34 PM EDT strict-source correction attempt} Curtis temporarily promoted a corrected actual-source crop for the same 2026-05-03 Scherzo-Tarantelle phrase. The detected sequence was \texttt{A\#4 D5 C5 A\#4 D5}; the displayed spelling was \texttt{Bb4 D5 C5 Bb4 D5}; the score-side sequence was the same octave-aware note sequence; and the score panel used a non-generated crop from the local Scherzo-Tarantelle solo part. This attempt is now superseded by Alan's later review: the phrase is rejected and cannot count as an accepted measure-level match, exact score-aligned window, score-coordinate heat-map fragment, score-visual lock, actual-source score lock, verified source snippet, or accepted five-note phrase.

\paragraph{2026-05-15 08:37 PM EDT fast-transition candidate update} The next blocker was not another UI issue: the broader pYIN/onset stream in the 2026-05-03 Scherzo-Tarantelle window still collapses large parts of fast playing into repeated D/B-flat/G material, so a longer source-score phrase cannot be honestly promoted. Curtis now adds a hidden YIN frame-transition trace to every transcription run. This trace follows fast frame-level pitch changes and stores them as \texttt{scoreMatchCandidateNotes}; it is never rendered as accepted notation by itself. The daily matcher can search those hidden candidates against parsed source-score notes, but accepted score evidence still requires exact MIDI equality, exact written score spelling, visible source-score crop agreement, local audio/video, and no uncertain octave-corrected detected notes. Regression tests now prove that hidden transition candidates can feed the strict score matcher while uncertain octave-corrected candidates are rejected. A local check on the 5/3 \texttt{Njh8\_zq9\_DM-9465} clip produced 174 hidden fast-transition candidate notes with 8 pitch classes. After the five-note rejection, those candidates remain search data only and do not create accepted evidence. This is implementation progress toward long phrases, not a solved long phrase.

\paragraph{2026-05-15 08:48 PM EDT notation-renderer correction} The web notation renderer now moves the treble-clef baseline upward and re-anchors flat key-signature glyphs so B-flat and E-flat render on the intended staff positions. This fixes a display-layer source of visible wrongness, but it does not count as transcription progress unless the paired audio and source-score sequence pass the strict acceptance gate.

\paragraph{2026-05-15 09:07 PM EDT targeted-transcription update} Live v11 reprocessing exposed a workflow bottleneck: the public transcribe endpoint could only walk the stale media queue in fixed eight-sample passes, so a known evidence clip could remain stale while unrelated samples were regenerated first. The endpoint now accepts a bounded \texttt{sampleId} query parameter and an optional bounded \texttt{limit}, allowing a specific local media sample to be regenerated under the current pipeline without bypassing violin-positive checks, hidden-candidate gating, uncertainty rejection, or source-score acceptance rules. This is required for faster long-phrase iteration because failures can be retested directly instead of waiting through broad queue churn.

\paragraph{2026-05-15 09:20 PM EDT live targeted-reprocess verification} The deployed endpoint was run against \texttt{Njh8\_zq9\_DM-9465}. The live tracker then reported 67.375 percent exact weighted progress, 1380 hidden transition-trace candidates, 1 accepted long-phrase counter item, 1 accepted measure-match counter item, and 1 exact-source visible match group on 2026-05-03. Browser verification of the public site showed one visible group: \texttt{Bb D C Bb D} in the score panel, local video/audio from \texttt{2:37:51--2:37:52}, and transcription \texttt{Bb4 D5 C5 Bb4 D5 / 1.1s}. Alan's later review rejected that group, so the retained progress is the targeted endpoint and hidden-candidate instrumentation, not the accepted phrase.

\paragraph{2026-05-15 09:24 PM EDT post-deploy state check} A post-deploy API read showed that live background processing had advanced the implementation tracker to 68.125 percent exact weighted progress and 2476 hidden transition-trace candidates. That one-group state is now superseded by the 10:05 PM EDT rejection lock. The extra hidden candidates improve search coverage, but they do not count as solved transcription until a longer exact source-score run is verified and reviewed.

\paragraph{2026-05-15 10:05 PM EDT five-note rejection lock} Alan's latest review rejected the five-note Scherzo-Tarantelle phrase itself. Curtis now demotes \texttt{A\#4 D5 C5 A\#4 D5} / \texttt{Bb4 D5 C5 Bb4 D5} from accepted evidence back to candidate-only material. Regression tests now require this exact lane to have zero verified score location, zero verified source snippet, no accepted measure, no accepted long phrase, no score-coordinate heat-map fragment, and no visible accepted score panel. The only items Alan had called usable remain short audio-only D/A detections; the A4/A5 score anchors, D--B-flat--G--D score sequence, and five-note Scherzo phrase are rejected until new independent review proves otherwise.

\paragraph{2026-05-15 10:54 PM EDT second-pass audio-gate update} The repeated loop was that fast transition notes and loose reference-audio overlaps could still look like progress before the notes were independently confirmed by audio. A local audit of the rejected \texttt{Njh8\_zq9\_DM-9465} clip read the candidate closer to \texttt{Bb4 D5 Bb4 D5}; the displayed middle \texttt{C5} did not survive the second pass. Curtis now preserves \texttt{audioAgreement}, \texttt{agreementSources}, \texttt{agreementSourceCount}, and detector provenance in daily detected-note series. Hidden \texttt{scoreMatchCandidateNotes} are filtered before score search unless each note has audio agreement, at least one independent agreement source, spectral/onset participation, confidence at or above 0.70, and enough duration to avoid transition glitches. Symbolic score matches and reference-sequence matches now require the same per-note audio gate, and reference-audio pitch-trace matches need at least four notes and three distinct pitch classes. Regression tests now lock the rejected Scherzo phrase when the \texttt{C5} fails audio agreement, reject unagreed fast candidate notes, reject uncertain octave-corrected candidates, and stop weak two-note reference-audio coincidences from counting as phrase progress. This does not solve long-phrase transcription; it removes a false-positive lane that was making Curtis go in circles.

\paragraph{2026-05-15 11:17 PM EDT candidate-isolation update} Curtis now separates source-verification candidates from displayable score matches in the daily-record API. Pitch-sequence candidates can still feed the truth workbench and source-verification target queue, but they no longer populate the visible \texttt{matchGroups} bucket, no longer mark a day as reference-linked, and no longer create heat-map fragments unless the exact accepted score-evidence gate passes. A new regression test proves that a locally playable Wieniawski source-sequence candidate can be queued with \texttt{candidateMatchGroups} while \texttt{matchGroups}, accepted score links, and heat-map fragments remain empty. This prevents another circular handoff where Alan has to inspect a displayed ``match'' and discover that it was only a candidate.

\paragraph{2026-05-15 11:32 PM EDT live candidate-isolation verification} The deployed API reported 46.125 percent exact weighted implementation progress, 0 accepted long phrases, 0 accepted measures, 0 exact score-aligned windows, 0 actual-source score snippets, and 0 score-coordinate heat-map fragments. The 2026-05-03 daily record reported \texttt{source\_verification\_pending}, 0 visible \texttt{matchGroups}, 6 queued \texttt{candidateMatchGroups}, 0 score links, and 0 reference links. The top queued source target is \texttt{A D D F\# D}; it remains rejected as accepted score evidence because it does not have a contiguous exact-MIDI run in the corrected source MusicXML.

\paragraph{2026-05-16 source-position and exact-MIDI update} The verified symbolic score map itself had a staff-position error. Rechecking the local IMSLP solo-violin page at high resolution and cross-checking the Wieniawski Society bars 5--9 example showed that the opening violin entrance is \texttt{A5 G5 F5 A5 G5 F5 A5 G\#5 F5}, not the prior \texttt{D6 C6 Bb5 D6 C6 Bb5 D6 C\#6 B5} map. The MusicXML now stores those nine actual source notes from measures 5--7. The matcher still uses pitch-class overlap for candidate search, but accepted symbolic score matches and source-verification targets now require exact detected MIDI equality, per-note audio agreement, and the same actual-source crop and truth-evidence gates. Regression tests now cover the corrected opening map, a rejected partial exact-MIDI target whose best overlap is only \texttt{A5 G5}, and a positive audio-agreed exact-MIDI phrase \texttt{G5 F5 A5 G\#5 F5} with MIDI \texttt{79 77 81 80 77}. This still does not accept a live long phrase; it removes the path that let a wrong-staff-position source map or pitch-letter-only target look solved.

\paragraph{2026-05-16 exact-MIDI truth-manifest update} The exact-MIDI rule is now a standalone implementation asset instead of only an inline scanner condition. The new \texttt{curtis-long-phrase-gold.json} manifest verifies one source-only Wieniawski positive phrase, \texttt{G5 F5 A5 G\#5 F5}, and blocks three known false phrase families: the old \texttt{D6 C6 Bb5} source map, the rejected \texttt{Bb4 D5 C5 Bb4 D5} lane, and the current queued \texttt{G5 G5 A\#5 A\#5 A\#5 A\#5 A5} target. The shared gate now reports contiguous exact-MIDI overlap, ordered overlap, exact source sequence, and audio-agreement state. The site exposes the source-positive and blocked-regression counts in the compact solve-status panel. This is not a solved live phrase; it is the regression harness needed so the next phrase attempt cannot reuse a known wrong score map or a pitch-letter-only match.

\paragraph{2026-05-17 Staff 4 source-lane extension} The next source-verification step extends the current accepted Staff 4 lane outward instead of starting a new random score match. The verified MusicXML source now carries the adjacent \texttt{Eb5 C5} notes to the right of the accepted \texttt{Eb5 Eb5 C5 Eb5 Eb5} window, creating a seven-note source lane \texttt{Eb5 Eb5 C5 Eb5 Eb5 Eb5 C5}. The first implementation state kept that extension source-only because the available paired audio did not yet pass the adjacent-note gate. The later right-1 promotion supersedes that blocker for the sixth note only: \texttt{Eb5 Eb5 C5 Eb5 Eb5 Eb5} is now accepted source/audio truth, while the right-2 seven-note lane still needs its own audit and truth acceptance. Regression tests require this separation so adjacent source growth cannot silently turn a partial prefix into a longer displayed phrase match.

\paragraph{2026-05-17 Staff 4 phrase-audit packet update} Curtis now has a dedicated generator for Staff 4 expansion decisions instead of retrying random score matches. The packet is keyed to the current source/audio lane, extracts the bounded local audio clip, attempts a matching local video clip, preserves stored detected notes, reruns pYIN, YIN, and spectral-peak pitch checks over the same audio, records onset times, writes a pitch-trace SVG, writes a spectrogram SVG, and stores the packet JSON under the runtime audit route. The first blocker packet preserved the \texttt{Eb5} versus \texttt{D5} failure. The later right-1 packet can accept the six-note lane only because exact MIDI, source crop, paired audio, and truth state all agree. The next packet must audit the right-2 seven-note lane independently.

\paragraph{2026-05-17 Staff 4 audit-rejection and next-window update} The original right-2 Staff 4 audit result is a rejected regression case in the long-phrase truth manifest. The rejected case binds the source target \texttt{Eb5 Eb5 C5 Eb5 Eb5 Eb5 C5} to the observed paired audio run \texttt{D\#5 D\#5 C5 D\#5 D\#5 D5 D\#5}; because the sixth note is \texttt{D5} instead of source \texttt{Eb5}, that exact expansion cannot reappear as accepted evidence. The expansion harness then tested the smaller right-1 window \texttt{Eb5 Eb5 C5 Eb5 Eb5 Eb5}; a later exact audio candidate for that smaller window now passes and is accepted. The rejected right-2 case remains blocked, and the next right-2 attempt must be a fresh audit rather than a reuse of the old failed run.

\paragraph{2026-05-17 Staff 4 adjacent-mining update} Curtis now runs a dedicated mining pass over stored May 3 detected audio-note windows for the exact Staff 4 right-1 MIDI sequence \texttt{75 75 72 75 75 75} and the longer right-2 sequence \texttt{75 75 72 75 75 75 72}. This is separate from score display and cannot accept evidence by itself. It reports exact audio-agreed candidates when they exist, otherwise it reports the nearest stored miss and its exact-prefix length. Regression tests cover both outcomes: the current bad run remains \texttt{not\_found} because it has \texttt{D5} where the source requires \texttt{Eb5}, while a synthetic raw detected run with \texttt{D\#5 D\#5 C5 D\#5 D\#5 D\#5 C5} becomes an exact mining candidate ready for audit.

\paragraph{2026-05-17 03:23 AM EDT Staff 4 source-audio rescanner update} Curtis now searches beyond the note windows already stored in the daily record. Starting from the accepted Staff 4 source/audio anchor, the backend locates the original local media sample, extracts a padded source-audio window around the accepted clip, reruns the note detector on that fresh audio, shifts the detected event times back into the source-video timeline, and adds those runs to both the anchored-expansion harness and the adjacent-mining pass. The cache key includes the media path, file size, modification time, source window, and transcription-pipeline version, so stale detector output is not silently reused after source media or extractor changes. This still cannot promote a longer phrase by itself: a rescanned run can only become a review candidate when its exact MIDI sequence matches the verified Staff 4 source lane, and it still needs paired-audio review, source crop agreement, local video/audio, and truth acceptance before it displays as a score match. The regression suite now proves that an exact \texttt{75 75 72 75 75 75 72} phrase found only in a fresh source-audio rescan enters \texttt{ready\_for\_truth\_review} and the mining pool, while failed or missing source media remains blocked instead of becoming visible evidence.

\paragraph{2026-05-17 Staff 4 anchor-persistence update} The accepted Staff 4 lane no longer depends on the current daily display retaining an accepted \texttt{matchGroups} card. Curtis now derives the anchor from the long-phrase truth manifest when the daily record has no visible accepted match group: source id \texttt{wieniawski-staff4-eb-eb-c-eb-eb-symbolic-v1}, source MIDI \texttt{75 75 72 75 75}, practice day \texttt{2026-05-03}, sample \texttt{Njh8\_zq9\_DM-8835}, source window \texttt{*8835-8925}, and local phrase timing. The source-audio rescanner, expansion harness, and adjacent-mining pass all call the shared anchor resolver and deduplicate the manifest anchor against an existing visible accepted anchor. A regression test now reproduces the live failure directly: the May 3 record has zero visible \texttt{matchGroups}, yet source-audio rescan still runs, phrase expansion still has an anchor, adjacent mining still searches the exact Staff 4 MIDI lane, and \texttt{no\_staff4\_anchor} does not return. This protects the fixed-lane workflow; it does not promote any new phrase.

\paragraph{2026-05-18 Staff 4 full-phrase audit correction} The audit route now distinguishes a true failed adjacent note from a current exact-MIDI phrase that is merely waiting for source/truth acceptance. Before this correction, the seven-note target \texttt{Eb5 Eb5 C5 Eb5 Eb5 Eb5 C5} could carry a full exact stored audio run while the packet generator still followed stale first-failure metadata and clipped only the first note for review. The \texttt{staff4\_phrase\_audit\_v3} packet ignores stale failed-note metadata whenever the current expansion already has matching target MIDI, audio MIDI, per-note audio agreement, and a source crop. The resulting packet is a full-phrase truth-review packet with local timing \texttt{20.225--23.280}, not a one-note blocker packet. A later 2026-05-18 truth promotion accepted that packet. That closed the former \texttt{pending\_source\_lock} blocker; the next source extent was then extended separately to an unaccepted eighth-note \texttt{A4} target.

\paragraph{2026-05-18 Staff 4 eighth-note source extension and notation pass} The next Staff 4 source-map step adds the immediate right-adjacent \texttt{A4} from the actual IMSLP source crop, making the verified source-only target \texttt{Eb5 Eb5 C5 Eb5 Eb5 Eb5 C5 A4} with MIDI \texttt{75 75 72 75 75 75 72 69}. Curtis bumped the source-audio rescanner cache to \texttt{staff4\_source\_audio\_rescan\_v8} and raised the adjacent guided target limit to eight notes so the system can test the new eighth note instead of stopping at source extent. The gate did not accept the expansion: local completion reports \texttt{blocked\_no\_audio\_candidate}, zero covered audio-note runs for the eight-note range, and \texttt{blocked\_media\_missing} for audit artifact generation because the Staff 4 sample is not present in runtime media storage. The notation renderer was also tightened: treble clef, noteheads, sharps, and flats now use Bravura/SMuFL music glyphs; B-flat/E-flat key-signature marks are locked to B and E staff positions; the notation viewBox has high-ledger headroom so upper notes do not clip; and regression tests prevent rough ellipse noteheads or generic clef glyphs from returning.

\paragraph{2026-05-18 03:35 AM EDT local-font notation-metrics update} Curtis now vendors the Bravura music font locally and uses the local file before any CDN fallback, so accepted notation does not depend on a remote font race. The front-end renderer has one shared notation metric set for staff line extent, note start and end positions, ledger-line width, stem offset, local accidental x/y offset, accidental viewport limits, and key-signature spacing. The visual change fixes the remaining hand-tuned accidental drift: sharps and flats now sit on the same vertical staff coordinate as the notehead they modify, first-note accidentals have enough space after the clef, B-flat/E-flat key signatures remain locked to the B and E staff positions, and browser QA verified SMuFL glyphs, local font loading, no generic clef fallback, and no ellipse noteheads.

\paragraph{2026-05-18 03:54 AM EDT notation-symbol spacing update} Curtis tightened the notation renderer a second time after review showed that the music glyphs were still too cramped and text-like. First-note spacing moved farther right, local accidentals moved farther left of the notehead, noteheads increased to 38 px, flats increased to 38 px, sharps increased to 36 px, ledger lines narrowed to match the larger notehead, and stem offsets were reduced so stems attach closer to the notehead edge. The backend symbolic-score SVG renderer now uses the same stem and accidental-size logic. Regression tests lock the first-note A-sharp clearance, the high D-sharp vertical alignment, the B-flat/E-flat key-signature positions, local Bravura font use, and the absence of rough ellipse noteheads or generic clef fallback. This is a visual notation-quality pass only; it does not change which transcription or score evidence is accepted.

\paragraph{2026-05-18 12:41 PM EDT engraved-notation scale correction} Browser QA showed that the accepted ABC engraving could still be squeezed by compact match-row layout, making professional glyphs appear too small even though the clef, notes, sharps, and flats came from the vendored engraving path. Curtis now renders accepted ABC at scale 1.05, uses a shorter staff width for stronger music proportions inside each snippet, keeps the hydrated notation lane at a 520 px minimum width, and gives narrow screens a thin internal horizontal scrollbar. Regression tests lock the notation minimum width and scrollbar rule so responsive layout cannot shrink noteheads, accidentals, or clef placement back into a cramped display. This is a display-quality correction only; it does not accept any new transcription or score evidence.

\paragraph{2026-05-17 12:56 PM EDT Staff 4 tiled-rescan update} Curtis now widens that source-audio search from one padded clip to a bounded tiled plan over the same source sample. For the accepted \texttt{*8835-8925} Staff 4 anchor, the backend generates the anchor-core tile plus four overlapping rightward tiles, each no longer than 18 seconds. Each tile has its own cache key, audio extract, detector run, shifted source-window timing, and merged audio-note run. Regression tests now require five Staff 4 rescan windows and prove that exact MIDI found only in those tiled source-audio runs still enters the expansion/mining pool without becoming visible accepted evidence.

\paragraph{2026-05-17 01:09 PM EDT Staff 4 anchor-reproduction update} Curtis now checks whether the fresh current-detector source-audio rescan can reproduce the accepted Staff 4 anchor before treating the widened scan as usable phrase-extension evidence. The live tiled pass produced five source-window tiles, four note runs, and 286 primary-plus-candidate note events, but adjacent mining still found zero exact Staff 4 right-extension candidates after 1547 stored/rescanned windows. The same pass also did not reproduce the accepted five-note anchor \texttt{75 75 72 75 75}. The tracker now records \texttt{not\_reproduced} instead of implying that anchor persistence equals fresh detector agreement. This changes the next implementation step: calibrate segmentation or accept a new current-detector anchor before expanding the Staff 4 phrase.

\paragraph{2026-05-16 09:07 PM EDT gold-review-loop update} Curtis now has a real review loop instead of only hidden candidates and scattered correction endpoints. The backend builds a queue from detected audio-note series and candidate score groups, choosing reviewable note windows that prefer audio-agreed, spectrally supported, non-repeated phrase material over repeated-note sludge. The new \texttt{/api/curtis/gold-review} route exposes accepted, rejected, pending, queue, score-ready, and evidence-ready counts. The \texttt{/api/curtis/gold-review/items} write path accepts pending, accepted, or rejected labels and mirrors accepted/rejected labels into the truth workbench. Audio-only accepted labels become gold audio truth; accepted score labels are refused unless the accepted audio notes and score notes match exactly by note name, octave, and MIDI sequence and include a source score coordinate. The site now includes a compact review surface with local video, audio, detected notation, editable accepted notes, optional score notes/location, and accept/reject actions. The solve-status tracker includes gold-label and queue counts. This still leaves live long phrases at zero; it creates the review data path required to make future phrase work accumulate instead of cycling through unverifiable candidates.

\paragraph{2026-05-16 09:36 PM EDT playable-window guard update} Live browser verification found a real media failure in the new review loop: one queued gold item pointed at a local audio extraction window longer than the playback endpoint permits. Curtis now rejects or shortens review-note windows before they enter the queue so every published gold candidate has a playable local clip route. A regression test covers a long-spaced detected series and proves the queued candidate stays within the endpoint's playback limit. This is not musical transcription progress by itself; it is an evidence-loop guard so the review queue cannot ask Alan to label a broken clip.

\paragraph{2026-05-20 score-transcription training split} Curtis keeps \texttt{transcription-alan} visible separately from \texttt{score-transcription}. The backend returns \texttt{audioQueue} previews for actual paired audio review and exposes a \texttt{score-transcription} lane even when the strict source gate has zero safe samples. The requested-repertoire source catalog remains visible as real source material, but generated score-transcription candidates are quarantined from the review queue unless the visible source crop carries a verified note sequence from that same crop. This fixes the false-pairing failure where Curtis showed a real IMSLP crop beside a generated phrase that had not been extracted from that crop. Future \texttt{score-transcription} packets can still test exact visual notation copying--pitches, octaves, accidentals, key signature, durations, rests, beams, tuplets, stem direction, notehead shape, spacing, slurs/ties, and source range--but only after source extraction proves that the displayed copy target is the same music as the displayed original source. \texttt{score-transcription} labels train score reading only; they do not become accepted audio-score evidence unless later audio, transcription, continuity, exact MIDI, source crop, and truth review gates all agree.

\subsection{Definition of solved}

Long-phrase transcription is solved for Curtis only when the system can take Alan's long-form practice archive and produce evidence groups that a musician can trust without manually checking every note. The complete goal is not merely a model that prints plausible notation. The solved system must satisfy all of the following conditions.

\begin{enumerate}
    \item \textbf{Full archive accounting.} Every practice video from the strict practice ledger is assigned to the correct practice day, has exact uploaded duration, and has a measured or explicitly pending active-practice state.
    \item \textbf{Active violin time.} Total practice time means intervals where Alan is actually playing violin. Talking, tuning, setup, walking, camera adjustment, silence, scrolling, and other non-playing intervals are excluded.
    \item \textbf{Local evidence playback.} Every accepted transcription snippet has a playable local video clip and audio clip in the Curtis site. Accepted evidence does not rely on a YouTube embed.
    \item \textbf{Accurate note events.} The transcribed notes match the paired audio, including repeated notes, fast scalar runs, arpeggios, shifts, string crossings, rests, and obvious restarts. A fast passage must not collapse into a false repeated pitch stream such as \texttt{D D D D}.
    \item \textbf{Rhythm and timing.} The system estimates onsets, durations, rests, tempo changes, slow-practice stretches, restarts, and repeated attempts. When rhythm is uncertain, the notation must mark uncertainty rather than invent a clean answer.
    \item \textbf{Professional notation.} Displayed notation uses a real notation engine or a notation-quality rendering path. Treble clef placement, staff scaling, note placement, stems, beams, rests, spacing, key signature, and measure layout must look like real sheet music, not a hand-drawn mockup.
    \item \textbf{Score-source truth.} A score snippet can appear only when the score-side notes are source-confirmed. Scanned images without a verified note map are not enough. A displayed A must point to an actual A in the score.
    \item \textbf{Score alignment.} For known repertoire, the played note sequence must align to a specific location in the score. The match can ignore rhythm during the first pitch-class search, but accepted score display needs exact score-note identity and a stable location.
    \item \textbf{Score-free exercise handling.} If Alan invents an exercise or practices a technique pattern that is not in the score, Curtis should still transcribe and group the pattern. It should not force the pattern into Scherzo-Tarantelle, Haydn, or another piece.
    \item \textbf{Repeat grouping.} Repeated attempts of the same passage are grouped as attempts, loops, restarts, or slow/faster passes. The UI should not dump the same measures repeatedly.
    \item \textbf{Heat maps.} Practice density, repetition density, problem density, and improvement layers appear on actual score coordinates when a score match exists, or on the generated pattern/transcription when the material is score-free.
    \item \textbf{Curtis-level observations.} Observations are specific and evidence-tied: what happened, where it happened, how often it happened, whether it improved, and why the issue blocks Curtis-level readiness. Generic advice is not accepted.
    \item \textbf{Compact interface.} Expanded daily entries stay compact. The useful unit is the grouped score--transcription--audio--video pair, not a large dashboard or a long explanatory panel.
    \item \textbf{Regression protection.} Wrong score-note boxes, failed transcriptions, unmatched notation, stale piece labels, and non-playing practice-time credit cannot silently reappear after a later change.
\end{enumerate}

\subsection{Acceptance gates}

The following gates define whether progress is real. If a change fails one of these gates, it is not counted as solved, even if the UI looks better.

\begin{description}[leftmargin=1.85em,style=nextline]
    \item[Archive gate.] The strict practice ledger must list every video from \texttt{violin 1} forward that is either violin-numbered or date-titled as practice. Same-day videos must merge into one daily record. Each record stores source video id, title, uploaded date, inferred practice date, uploaded duration, and processing coverage.
    \item[Practice-time gate.] Active violin time is the sum of detected playing intervals, not the uploaded duration. The API must expose measured active time, checked video time, unmeasured video time, and estimate status separately. The UI label should use total practice time for active violin time and uploaded video for total source duration.
    \item[Evidence-media gate.] A transcription line is accepted only if the exact audio and video window are available locally, playable in the same Curtis page, and time-linked to the source record.
    \item[Note gate.] A displayed note must agree with the paired audio. The note label, octave where claimed, and staff position must be internally consistent. If the audio says A4, the rendered staff note must be A4, not F, G, B, or a visually adjacent note.
    \item[Score gate.] An accepted score snippet must come from the actual piece source: a verified source-score crop, verified MusicXML engraving export, or manually verified note-coordinate map for the same reference range. Generated notation may show the transcription, but it cannot appear as the accepted score snippet.
    \item[Single-note gate.] A one-note score display is allowed only as a calibration anchor, and only if both sides explicitly show the same pitch name and the score-side coordinate is verified. It must not be used as proof of a long passage.
    \item[Phrase gate.] A real phrase match requires a contiguous pitch-class sequence in the score or pattern source. Rhythm can be ignored for early matching, but accepted phrase display must keep the exact played sequence, the exact score sequence, and any omitted or uncertain notes.
    \item[Notation gate.] Hand-built notation SVG is not the long-term display path. The production path must use a professional renderer or renderer-equivalent layout with verified clef, key, staff, note, rest, spacing, and responsive behavior.
    \item[Exercise gate.] If no score match exists and the material forms a repeated exercise, the system must label it as an exercise or pattern, preserve the generated transcription, and avoid adding a fake repertoire entry.
    \item[Observation gate.] A Curtis-level observation must cite the clip, transcription event, score measure or pattern segment, issue type, repeat count, direction of change, and readiness consequence.
    \item[UI gate.] The day row stays chronological, the expanded state is compact, and only grouped evidence appears by default. Raw untrusted transcription stays hidden or goes to the transcription-run PDF with clear status.
    \item[Paper gate.] Any meaningful progress changes this section, the system-state table, and the limitation text in the same work cycle.
\end{description}

\subsection{Required source material}

The implementation cannot be completed by code alone. It needs a source bundle that makes the audio, score, and corrections checkable. The Wieniawski Scherzo-Tarantelle solo PDF is now present locally as the first source-score asset, and its first verified symbolic map covers the nine-note opening phrase only.

\begin{enumerate}
    \item \textbf{Raw or locally captured practice media.} Curtis needs stable local access to audio/video windows from the YouTube archive or owner-provided files. Metadata alone cannot produce practice time, clips, or transcription.
    \item \textbf{Source-confirmed piece names by day.} Alan can supply piece names when automatic detection is uncertain. These labels should seed matching but never become score evidence by themselves.
    \item \textbf{Symbolic scores for known repertoire.} For each piece, Curtis needs MusicXML, MEI, a verified MuseScore export, or an equivalent symbolic note map with measure numbers, pitches, durations, voices, and staff positions.
    \item \textbf{Score images tied to symbolic coordinates.} To show actual score snippets, the image/PDF rendering must be linked to the same symbolic note ids or verified coordinate boxes.
    \item \textbf{Manual corrections.} Alan or a reviewer must be able to correct a note, octave, onset, rhythm, score location, or false match. These corrections become training and regression data.
    \item \textbf{Calibration recordings.} Labeled scales, arpeggios, open strings, repeated notes, shifts, slow exercises, and fast exercises help tune pitch detection before attempting long phrases.
    \item \textbf{Gold benchmark clips.} A small set of manually checked audio--transcription--score pairs is required before the app can claim improvement. The benchmark must include easy anchors, repeated notes, fast runs, arpeggios, and real Scherzo-Tarantelle fragments.
\end{enumerate}

\subsection{Implementation plan}

\subsubsection*{Phase 0: Make the paper the tracker}

\textbf{Status: complete in this revision.} The first required step is to make the plan durable in the Curtis paper PDF. The paper now records the definition of solved, the current baseline, the gates, the phase plan, and the update protocol. Future changes are not complete unless the paper is checked or updated.

\subsubsection*{Phase 1: Build the canonical evidence ledger}

The backend needs a stricter data model for practice evidence. The current modules already include the daily-record builder, media sampler, transcription engine, symbolic-score parser, score-asset layer, reference corpus, and correction layer. The next version should make the following entities explicit:

\begin{itemize}
    \item \texttt{source\_video}: immutable source id, title, URL, uploaded time, duration, practice-date basis, and ledger inclusion basis.
    \item \texttt{practice\_day}: date, list of source videos, uploaded duration, active-practice coverage, piece labels, and status.
    \item \texttt{active\_interval}: source id, start, end, confidence, detector version, reason for inclusion, and reason for exclusion when withheld.
    \item \texttt{media\_clip}: local video path, local audio path, waveform path, source interval, duration, and playable route.
    \item \texttt{note\_event}: onset, offset, pitch class, octave, cents deviation, confidence, detector sources, uncertainty reason, and correction state.
    \item \texttt{score\_note}: piece id, movement, measure, voice, pitch class, octave, duration, note id, and rendered coordinates.
    \item \texttt{match\_candidate}: played note ids, score note ids or exercise-pattern ids, match score, mismatch list, and display eligibility.
    \item \texttt{accepted\_evidence}: only the groups that pass the gates; these are the only groups the main UI can show.
    \item \texttt{rejected\_evidence}: wrong notes, wrong score boxes, false labels, failed transcriptions, and model guesses that must not reappear.
\end{itemize}

The acceptance condition for this phase is that a wrong A-to-B or A-to-G score crop can be stored as rejected evidence and cannot be selected by the display builder again.

\subsubsection*{Phase 2: Complete the full active-practice scan}

The practice-time problem is separate from transcription. Curtis should first scan the full 276 h 23 min uploaded practice ledger for violin-playing intervals with a low-cost pass. The pass should use audio energy, violin-range voiced pitch, onset density, spectral shape, and visual fallback only where necessary. It should not require full transcription.

The fast pass should walk every practice video in chronological order and store approximate active intervals. A second exact pass should refine boundaries around detected regions. The UI should show three numbers for every day: uploaded video time, checked video time, and total practice time. While scanning is partial, the estimate stays labeled as estimated. When every video has been scanned, the estimate should be replaced by measured total practice time.

Acceptance for this phase:
\begin{itemize}
    \item every strict-ledger video has a coverage record;
    \item every day has uploaded duration and active-practice status;
    \item non-playing windows are excluded;
    \item active-practice total can be recomputed from stored intervals;
    \item the paper table is updated with checked duration, active duration, unmeasured duration, and estimate status.
\end{itemize}

\subsubsection*{Phase 3: Generate stable local evidence clips}

Every accepted interval and every candidate transcription window needs local media. Curtis should generate short MP4 clips and audio clips with stable IDs derived from source video id and time range. The browser should play the local clip and audio without YouTube embedding. The audio used for transcription must be the same audio the user can play.

Acceptance for this phase:
\begin{itemize}
    \item each accepted transcription event links to local video and local audio;
    \item the site can play the clip in the same window;
    \item source time, local file, and rendered transcription agree;
    \item dead local media routes fail the evidence gate rather than rendering a broken accepted group.
\end{itemize}

\subsubsection*{Phase 4: Build the verified score library}

The current weak link is score-side truth. Curtis must stop relying on guessed scanned-score boxes. For Wieniawski Scherzo-Tarantelle and any future piece, the score library needs symbolic notes and rendered coordinates.

The best path is:
\begin{enumerate}
    \item acquire a legal public-domain score image or PDF;
    \item obtain or create a symbolic MusicXML/MEI/MuseScore representation;
    \item render the symbolic score to image or SVG with stable note ids;
    \item map each rendered note to pitch, octave, measure, and coordinates;
    \item manually verify a representative set of anchors, especially A, D, G, B, fast runs, ledger lines, accidentals, and octave-displaced passages;
    \item reject any score source where the note map cannot distinguish adjacent notes reliably.
\end{enumerate}

Acceptance for this phase:
\begin{itemize}
    \item a score-side A in the UI is generated from an actual score-side A;
    \item the rendered crop includes enough context to be legible;
    \item the crop does not cut off clefs, notes, beams, stems, or measure context;
    \item single-note anchors are visually verified before display;
    \item phrase snippets use the exact matched symbolic notes, not a guessed image region.
\end{itemize}

\subsubsection*{Phase 5: Replace fragile note extraction with a multi-stage transcription engine}

The note detector needs to work on violin practice, not only on clean isolated notes. The engine should process active intervals in stages:

\begin{enumerate}
    \item normalize audio and remove long silence;
    \item isolate violin-dominant energy where possible;
    \item detect onsets and repeated attacks before pitch smoothing;
    \item run multiple pitch estimators over the same window;
    \item separate stable pitch, glissando/shift, vibrato, noise, and uncertain frames;
    \item group frames into note events by onset, pitch stability, and duration;
    \item preserve fast pitch changes instead of smoothing them into one repeated pitch;
    \item estimate rests and pauses between attempts;
    \item detect double stops or chords as multi-note events when monophonic assumptions fail;
    \item attach confidence and uncertainty to every event;
    \item store rejected alternatives so repeated failures can be debugged.
\end{enumerate}

Acceptance for this phase:
\begin{itemize}
    \item open strings and simple scales transcribe correctly;
    \item repeated notes remain repeated notes with distinct onsets;
    \item arpeggios become changing pitch sequences, not a repeated-note collapse;
    \item obvious fast runs preserve the contour and most pitch classes;
    \item uncertain events are labeled uncertain rather than shown as clean notation;
    \item every displayed note can be played from its paired audio.
\end{itemize}

\subsubsection*{Phase 6: Make transcription score-aware}

For known repertoire, Curtis should not transcribe in isolation and then guess the piece later. The user-supplied or detected piece name should provide a score prior. The aligner should search the score for pitch-class sequences that match the audio-derived note stream, allowing restarts, skipped notes, slow practice, repeated loops, tempo changes, and rhythm mismatch during early search.

The alignment should use multiple layers:
\begin{itemize}
    \item pitch-class subsequence search for easy anchors;
    \item interval-contour search when absolute pitch is uncertain;
    \item dynamic time warping for variable tempo;
    \item Viterbi-style path selection through the symbolic score for longer passages;
    \item restart and loop detection when the path jumps backward;
    \item score-free fallback when no score path is plausible.
\end{itemize}

Acceptance for this phase:
\begin{itemize}
    \item one-note anchors pass only when the score coordinate is verified;
    \item three-to-five-note sequences can be matched to exact score notes;
    \item longer phrases show played notes beside matched score notes;
    \item mismatches are visible rather than hidden;
    \item rhythm mismatch can be tolerated during matching but not during final notation labeling;
    \item accepted score matches survive a musician spot-check.
\end{itemize}

\subsubsection*{Phase 7: Support score-free technique exercises}

Alan may invent an exercise during practice. Curtis must treat this as first-class evidence. If the note stream is coherent but does not match the named score, the system should cluster repeated patterns, render generated notation, and label the material as a technique pattern or exercise.

Acceptance for this phase:
\begin{itemize}
    \item repeated made-up patterns group together;
    \item the generated transcription is linked to audio/video;
    \item no fake score snippet appears;
    \item the repertoire entry records technique-pattern evidence separately from piece evidence;
    \item heat maps can show pattern density without pretending there is a score location.
\end{itemize}

\subsubsection*{Phase 8: Add manual correction and training}

Fully automatic transcription will not become reliable without corrections. Curtis needs a correction UI that allows Alan to say: this note is A, this staff note is wrong, this score box points to B, this clip is non-playing, this rhythm is wrong, this is Scherzo-Tarantelle, this is a made-up exercise, or this score match is false.

\textbf{Status: partially implemented.} The gold review loop now covers \texttt{transcription-alan} acceptance/rejection, score-phrase acceptance/rejection from queued detected clips, and \texttt{score-transcription} acceptance/rejection from visible source snippets. The \texttt{score-transcription} lane is deliberately separate: it trains whether Curtis copied the original score notes correctly, but it does not create accepted playing evidence, score-linked phrase evidence, or solved transcription by itself. The loop remains intentionally narrow: it can create accepted audio truth, exact score-ready truth, and score-reading labels, but it does not yet provide detailed onset/rhythm editing, score-box coordinate editing, batch keyboard review, waveform overlays, or full measure-level correction.

Every correction should become data:
\begin{itemize}
    \item corrected note events become gold labels;
    \item rejected matches become negative examples;
    \item accepted matches become benchmark fixtures;
    \item corrected piece names update source labels without unlocking score snippets by themselves;
    \item corrections invalidate stale derived evidence.
\end{itemize}

Acceptance for this phase:
\begin{itemize}
    \item Alan can accept or reject queued note windows without editing code;
    \item accepted and rejected examples appear in the gold/truth set;
    \item the same failure cannot reappear after rerendering;
    \item the PDF tracker records the correction count and benchmark count.
\end{itemize}

\subsubsection*{Phase 9: Build the benchmark suite}

Curtis needs objective tests before claiming improvement. The benchmark should include:

\begin{itemize}
    \item open-string clips;
    \item G, D, A, and E major/minor scales where available;
    \item arpeggios;
    \item repeated single-note attacks;
    \item fast scalar runs;
    \item shifts with sliding or unstable pitch;
    \item Scherzo-Tarantelle fragments with verified score locations;
    \item Haydn fragments with verified score locations;
    \item score-free exercises;
    \item non-playing windows that must not count as practice time.
\end{itemize}

Metrics should include note precision, note recall, pitch-class accuracy, octave accuracy when claimed, onset error, duration error, false-note rate, repeated-note recall, phrase-match precision, phrase-match recall, score-coordinate accuracy, false accepted score match count, active-practice precision, active-practice recall, and UI evidence completeness.

Acceptance for this phase:
\begin{itemize}
    \item tests fail on the known A-to-wrong-score-note examples;
    \item tests fail on repeated-pitch collapse;
    \item tests fail when accepted notation lacks paired audio/video;
    \item tests fail when practice time includes non-playing footage;
    \item tests pass only when the displayed evidence is actually trustworthy.
\end{itemize}

\subsubsection*{Phase 10: Use professional notation rendering}

The current custom notation path is only a temporary evidence display. The solved system should render notation from MusicXML, MEI, or another structured notation representation through a professional renderer or renderer-equivalent layout. The renderer must handle treble clef placement, staff scale, ledger lines, rests, stems, beams, ties, slurs where known, key signatures, accidentals, note spacing, measure widths, and responsive crop behavior.

Acceptance for this phase:
\begin{itemize}
    \item treble clef is correctly centered around the G-line spiral;
    \item the staff note for A4 appears in the A4 position;
    \item key signature and accidentals follow the chosen notation policy;
    \item snippets look publication-quality at desktop and mobile widths;
    \item the same notation source can export to the transcription-run PDF.
\end{itemize}

\subsubsection*{Phase 11: Build heat maps from matched coordinates}

Heat maps should be computed from actual practice behavior. For score-matched passages, each matched score note or measure receives density values. For score-free exercises, each pattern segment receives density values.

The layers are:
\begin{itemize}
    \item practice density: active time spent on the passage;
    \item repetition density: number of attempts or loops;
    \item problem density: restarts, stops, slowdowns, unstable notes, and failed attempts;
    \item improvement: later attempts cleaner than earlier attempts.
\end{itemize}

Acceptance for this phase:
\begin{itemize}
    \item a heat-map mark can be traced to clips and note events;
    \item a score-free exercise can still have a pattern heat map;
    \item heat maps do not appear as decorative overlays when no match exists;
    \item the day-level heat map and repertoire heat map use the same evidence.
\end{itemize}

\subsubsection*{Phase 12: Generate specific Curtis-level observations}

Observation generation should happen after transcription and alignment, not before. The observation engine should compare attempts of the same passage and classify issues as pitch, rhythm, bow, shift, tone, articulation, tempo, consistency, musical shape, or visible tension.

The output for each main passage should include:
\begin{itemize}
    \item likely passage or exercise pattern;
    \item evidence clip;
    \item transcription snippet;
    \item score snippet when verified;
    \item exact issue;
    \item number of attempts where it occurred;
    \item whether it improved, stayed the same, or got worse;
    \item why it blocks Curtis-level readiness.
\end{itemize}

Acceptance for this phase:
\begin{itemize}
    \item observations cite evidence;
    \item generic advice is absent;
    \item a day ends with one specific main blocker;
    \item rejected or uncertain transcription does not feed readiness claims.
\end{itemize}

\subsubsection*{Phase 13: Keep the UI compact}

The main page should show the paper card, analyzed practice days, and repertoire. Expanded daily entries should show one or a small number of grouped evidence pairs, not a full debug dashboard. The day marker should remain visible on the left so boundaries between days are clear. Videos should be small enough that a day entry fits on screen when possible.

Acceptance for this phase:
\begin{itemize}
    \item one daily group contains score, clip, audio, transcription, and the minimum useful status;
    \item raw text explanations are removed from the main view;
    \item chronological order is stable;
    \item desktop and mobile render without overlap, excessive height, or broken media controls;
    \item unresolved material is visible as state, not as a fake dashboard.
\end{itemize}

\subsubsection*{Phase 14: Automate paper updates}

The paper is now the progress tracker. Future meaningful Curtis changes should update:
\begin{itemize}
    \item live state numbers;
    \item solved gates;
    \item benchmark metrics;
    \item accepted score-match count;
    \item accepted transcription duration;
    \item full archive scan coverage;
    \item remaining blockers;
    \item limitations.
\end{itemize}

Acceptance for this phase is procedural: after every Curtis change that affects transcription, score matching, practice time, clips, heat maps, repertoire, observations, or evidence gates, the LaTeX paper is checked, rebuilt, and the root \texttt{paper.pdf} is refreshed.

\subsection{Milestones}

\begin{table}[htbp]
\centering
\scriptsize
\setlength{\tabcolsep}{4pt}
\renewcommand{\arraystretch}{1.1}
\caption{Milestone plan for moving from current evidence gates to solved long-phrase transcription.}
\label{tab:roadmap_milestones}
\begin{tabularx}{\textwidth}{>{\bfseries}p{0.12\textwidth}p{0.18\textwidth}X p{0.18\textwidth}}
\toprule
Milestone & Target & Acceptance evidence & Status \\
\midrule
M0 & Paper tracker & Roadmap is included in the Curtis paper PDF and becomes the update record. & Complete in this revision. \\
M1 & Practice time & All 276 h 23 min of strict-ledger video has active-practice coverage and per-day totals. & Active-scan route deployed; live state now has 331 active intervals from 106 sample results, 2 h 24 min checked video, 2 h 21 min measured active practice, and 250 queued windows; full chronological scan still pending. \\
M2 & Score truth & Scherzo-Tarantelle has verified symbolic notes and score coordinates. & Partial; verified symbolic source notes remain available, but the May 3 measure-16 visible lane is revoked. No score crop may display as accepted until the source notes, measure location, boxed crop, transcription, and paired audio agree. \\
M3 & One-note anchor & A displayed A from audio is matched to a verified score-side A, with no wrong-note crop. & Not accepted yet; audio-only D/A fragments are useful, but all A-specific score anchors remain rejected until the score-side notehead and octave are verified. \\
M4 & Short sequence & Three-to-five played notes align to exact score notes and render with local audio/video. & Blocked by evidence gate; previous Staff 4 short-sequence display is rejected after visible score/transcription mismatch, so the next accepted short sequence must start from a newly verified crop, audio, video, and transcription pair. \\
M5 & Long phrase & A 10--30 s passage transcription matches the paired audio and score location. & Pending. \\
M6 & Repeat grouping & Multiple attempts of the same passage group into loops, restarts, slow attempts, and faster attempts. & Pending. \\
M7 & Heat map & Practice density and problem density appear on verified score coordinates. & Pending. \\
M8 & Repertoire update & Piece entries update from accepted daily evidence and show active time, clips, snippets, and blockers. & Pending. \\
M9 & Readiness observation & Main blocker is generated from accepted evidence, not generic advice. & Pending. \\
M10 & Regression lock & Tests prevent failed notation, wrong score boxes, repeated rejected review candidates, and non-playing practice-time credit from returning. & Advanced; current tests block the user-rejected May 3 Staff 4 source ranges from source snippets, anchor selection, raw detected-series promotion, source-audio rescanning, and live display. Gold-review tests now also hide exact normalized MIDI patterns, same clip/score fingerprints, overlapping rejected clip windows, and already accepted source areas. Six rejected regressions are locked; no live Staff 4 accepted phrase is allowed until a fresh source crop is verified. \\
\bottomrule
\end{tabularx}
\end{table}

\subsection{Immediate next implementation}

After this PDF tracker, the first concrete implementation target is \textbf{full active-practice coverage plus benchmark/correction scaffolding}. This is the first useful code milestone because transcription cannot be trusted without knowing which source windows contain violin playing, and it cannot improve reliably without a gold correction loop. The first gold review loop is now implemented for queued audio/score phrase labels, but the larger benchmark suite and full archive coverage remain incomplete.

The 2026-05-13 11:57 AM EDT backend update completed the coverage and correction scaffolding portion:
\begin{itemize}
    \item added a canonical active-practice coverage ledger exposed through the review API and \texttt{/api/curtis/active-practice-coverage};
    \item separated uploaded video duration, checked sampled video duration, measured active practice time, candidate violin windows, unmeasured video time, and estimate state;
    \item counted checked non-playing sampled windows as checked video without adding them to practice time;
    \item added \texttt{/api/curtis/evidence-corrections} for note, score, match, piece, and active-interval corrections;
    \item added benchmark fixtures for rejected score-note and match corrections;
    \item added regression tests for checked non-playing windows, wrong score-note benchmark creation, and rejecting accepted score evidence when the displayed score note conflicts with the observed note.
\end{itemize}

The 2026-05-13 12:33 PM EDT backend update completed the active-scan route portion:
\begin{itemize}
    \item added a persistent \texttt{activePracticeScan} state record with active intervals, sample results, pending windows, and run history;
    \item scans stored local media samples only after the violin-presence gate is positive;
    \item extracts sound-active sub-intervals inside the source window and stores them as active-practice intervals;
    \item records negative violin-presence samples as checked without adding practice time;
    \item feeds persisted active intervals into per-day totals and the active-practice coverage ledger;
    \item exposes active-scan status and run endpoints so the deployed runtime can perform the pass without transcription;
    \item queues the next unchecked 90-second strict-ledger windows for owner media capture.
\end{itemize}

The 2026-05-13 12:59 PM EDT deployed scan update completed the first live run and exposed a practice-time accounting failure mode: sparse, captured samples can make an archive-wide estimate look too certain. The implementation now records active-scan status counts, caps active seconds by checked source coverage, and treats active-scan results as the practice-time source for scanned samples instead of letting noisy transcription durations override them.

The 2026-05-13 03:26 PM EDT queue update connected the active-scan queue to the owner-media sync path. The sync path now asks the deployed active-scan run endpoint for pending windows, keeps the response compact, and places those windows ahead of generic expansion or blind sampling. The pending queue remains the source of work for M1: each captured window must still be downloaded or browser-captured, locally classified for violin presence, uploaded, rescanned, and recorded before it becomes checked practice-time coverage.

The 2026-05-13 03:37 PM EDT owner-sync run proved that the local upload/rescan loop still works after the queue update. It uploaded one cached violin-positive 5/3 window and the deployed active scan converted it into one additional active-practice interval. Because cached violin-positive media is still prioritized ahead of new capture, this run did not consume the first chronological pending window. The next M1 step should either exhaust the cached positive backlog first or run a queue-only owner-sync mode so the pending list starting at \texttt{violin 1} advances.

The 2026-05-13 04:05 PM EDT update added that queue-only owner-sync mode and used it to consume the first chronological pending window. The \texttt{violin 1} \texttt{*0-90} window was captured, classified as non-playing/unclear, uploaded, and folded into the scan results as checked non-violin evidence. This is correct progress because it prevents setup or non-playing footage from being counted as violin practice. It does not advance transcription because no violin-positive musical material was present in that window. The UI now shows the current long-phrase transcription completion percentage and the complete done/open gate list at the bottom of the Curtis page so the remaining work is visible without turning daily entries into a large debug dashboard.

The 2026-05-13 05:44 PM EDT update advanced the same queue from \texttt{*90-180} through \texttt{*450-540}. Those windows produced active-practice evidence and moved checked coverage to 2 h 3 min. The timeout exposed a resume failure: cached active-scan pending windows could be captured but then skipped on the next queue-only run. The helper now includes matching cached pending-window files in queue-only order, including non-playing cached windows, so an interrupted scan can resume rather than forcing Alan or Codex to manage the cache manually. The coverage ledger also now treats unchecked transcription-only remnants as zero practice time; active practice cannot exceed checked video coverage.

The 2026-05-13 06:01 PM EDT update consumed the next chronological queue window, \texttt{*540-630}. That window was violin-positive and added four active intervals after the active-practice scan. The run succeeded without a timeout at batch size one, and the next pending window became \texttt{*630-720}. This kept the current work focused on M1 coverage rather than score matching or long-phrase notation.

The 2026-05-13 06:26 PM EDT update consumed the next chronological queue window, \texttt{*630-720}. That window was violin-positive and added five active intervals after the active-practice scan. The command wrapper timed out while the owner-sync process was still running, so the process was allowed to finish and then verified from live state instead of starting a duplicate run. The next pending window is now \texttt{*720-810}.

The 2026-05-13 07:55 PM EDT update consumed the next chronological queue window, \texttt{*720-810}. That window was violin-positive and added three active intervals after the active-practice scan. The owner-sync process completed cleanly with a longer timeout, so no duplicate capture or manual cache recovery was needed. The next pending window is now \texttt{*810-900}.

The 2026-05-13 08:41 PM EDT status-surface update made the Curtis page itself carry the implementation state instead of leaving the percentage and plan in the chat thread. The bottom tracker is now labeled as long-phrase transcription implementation status and is driven by the same backend gate object as the API. The visible panel reports 29 percent completion, 2 h 7 min checked video out of 213 h 23 min, 2 h measured active practice, 0 exact score windows, 0 accepted long phrases, 250 queued windows, and 250 active intervals. It also exposes a compact done list, an open-work list, an eight-phase implementation plan, and a collapsed detailed gate checklist. This is a presentation and tracking improvement, not a transcription breakthrough: score-truth verification, phrase-level note/rhythm extraction, exact score alignment, heat maps, and Curtis-level observations remain open.

The 2026-05-13 09:25 PM EDT implementation pass advanced the M1 scan and M10 benchmark loop. The queue-only sync consumed \texttt{violin 1} windows \texttt{*810-900} and \texttt{*900-990}; both became violin-positive source evidence and the deployed active-practice scan folded them into the checked ledger. Live state now reports 31 percent implementation completion, 2 h 10 min checked video out of 213 h 23 min, 2 h 3 min measured active practice, 92 sample results, 82 active-violin samples, 10 checked non-violin samples, 257 active intervals, and 250 queued windows. The correction ledger now contains three benchmark fixtures for real wrong-score-note failures: \texttt{A4}/B-crop, \texttt{A4}/\texttt{G4}-crop, and \texttt{A5}/\texttt{G5}-crop. Those cases are progress only because they prevent the same bad score matches from being counted as solved evidence. Curtis still has 0 accepted exact score windows and 0 accepted long phrases.

The 2026-05-13 10:05 PM EDT implementation pass continued the M1 scan and fixed a tracker presentation gap. The queue-only sync consumed \texttt{violin 1} windows \texttt{*990-1080} and \texttt{*1080-1170}; both became violin-positive source evidence and the deployed active-practice scan folded them into the checked ledger. Live state now reports 31.13 percent exact weighted implementation completion, 2 h 13 min checked video out of 213 h 23 min, 2 h 6 min measured active practice, 94 sample results, 84 active-violin samples, 10 checked non-violin samples, 269 active intervals, and 250 queued windows. The rounded public label still reads 31 percent, but the API and site now expose exact weighted progress so small M1 coverage gains do not disappear. Curtis still has 0 accepted exact score windows and 0 accepted long phrases.

The 2026-05-13 11:14 PM EDT implementation pass continued M1 and tightened the progress calculation itself. The queue-only sync consumed \texttt{violin 1} windows \texttt{*1170-1260} and \texttt{*1260-1350}; the first became violin-positive source evidence and the second became checked non-playing/unclear evidence. The deployed active-practice scan folded both into the coverage ledger. Live state now reports 96 sample results, 85 active-violin samples, 11 checked non-violin samples, 283 active intervals, 2 h 16 min checked video, 2 h 7 min measured active practice, and a 199 h 19 min provisional archive estimate. The tracker implementation now sums unrounded gate points before formatting the public label, so this small but real coverage gain moves the expected exact implementation label to 31.128 percent after deployment. Curtis still has 0 accepted exact score windows and 0 accepted long phrases.

The 2026-05-13 11:38 PM EDT implementation pass continued M1. The queue-only sync consumed \texttt{violin 1} windows \texttt{*1350-1440} and \texttt{*1440-1530}; both became violin-positive source evidence and the deployed active-practice scan folded them into the checked ledger. Live state now reports 31.131 percent exact weighted implementation completion, 2 h 19 min checked video out of 213 h 23 min, 2 h 10 min measured active practice, 98 sample results, 87 active-violin samples, 11 checked non-violin samples, 294 active intervals, and 250 queued windows. This is real practice-time coverage progress only. Curtis still has 0 accepted exact score windows and 0 accepted long phrases.

The 2026-05-13 11:55 PM EDT implementation pass continued M1. The queue-only sync consumed \texttt{violin 1} windows \texttt{*1530-1620} and \texttt{*1620-1710}; the first became checked non-playing/unclear evidence and the second became violin-positive source evidence. The deployed active-practice scan folded both into the checked ledger. Live state now reports 31.134 percent exact weighted implementation completion, 2 h 22 min checked video out of 213 h 23 min, 2 h 12 min measured active practice, 100 sample results, 88 active-violin samples, 12 checked non-violin samples, 300 active intervals, and 250 queued windows. This is real practice-time coverage progress only. Curtis still has 0 accepted exact score windows and 0 accepted long phrases.

The 2026-05-14 12:09 AM EDT implementation pass continued M1. The queue-only sync consumed \texttt{violin 1} windows \texttt{*1710-1800} and \texttt{*1800-1890}; both became violin-positive source evidence. The deployed active-practice scan folded both into the checked ledger. Live state now reports 31.136 percent exact weighted implementation completion, 2 h 25 min checked video out of 213 h 23 min, 2 h 15 min measured active practice, 102 sample results, 90 active-violin samples, 12 checked non-violin samples, 313 active intervals, and 250 queued windows. This is real practice-time coverage progress only. Curtis still has 0 accepted exact score windows and 0 accepted long phrases.

The 2026-05-14 12:35 AM EDT implementation pass continued M1. The queue-only sync consumed \texttt{violin 1} windows \texttt{*1890-1980} and \texttt{*1980-2070}; both became violin-positive source evidence. The deployed active-practice scan folded both into the checked ledger. Live state now reports 31.139 percent exact weighted implementation completion, 2 h 28 min checked video out of 213 h 23 min, 2 h 18 min measured active practice, 104 sample results, 92 active-violin samples, 12 checked non-violin samples, 322 active intervals, and 250 queued windows. This is real practice-time coverage progress only. Curtis still has 0 accepted exact score windows and 0 accepted long phrases.

The 2026-05-14 03:00 AM EDT implementation pass continued M1. The queue-only sync consumed \texttt{violin 1} windows \texttt{*2070-2160} and \texttt{*2160-2250}; both became violin-positive source evidence. The deployed active-practice scan folded both into the checked ledger. Live state now reports 31.142 percent exact weighted implementation completion, 2 h 31 min checked video out of 213 h 23 min, 2 h 21 min measured active practice, 106 sample results, 94 active-violin samples, 12 checked non-violin samples, 331 active intervals, and 250 queued windows. This is real practice-time coverage progress only. Curtis still has 0 accepted exact score windows and 0 accepted long phrases.

The 2026-05-14 05:11 AM EDT source-score acceleration pass advanced M2 and M4 instead of continuing only the slow archive scan. The implementation added a local MusicXML source excerpt for the Scherzo-Tarantelle solo-violin opening motif, wired it into the Wieniawski score target, and changed the symbolic matcher so this source target can accept a four-note measure-level run with at least three distinct pitch classes. The parser now preserves written flats such as B-flat for staff placement while still matching their pitch class to detected audio, which closes the failure mode where a normalized A-sharp could render on the wrong staff position. The score snippet renderer now receives the target key signature and renders a G-minor two-flat key signature before the matched notes. The completion tracker now separates \texttt{acceptedMeasureMatchCount} from \texttt{longPhraseAcceptedCount}, so Curtis can show one verified measure-level score/audio match without claiming solved long-phrase transcription. The regression suite increased to 103 tests and now covers source-backed Wieniawski symbolic notes, four-note source-measure matching, written-flat staff placement, accepted-measure counting, and the continued rejection of one-note visual crops. This pass does not solve long-phrase transcription; it creates the first narrow, test-covered lane from source score notes to a grouped score--audio--video--transcription match.

The 2026-05-16 10:30 PM EDT source-review pass completed the next narrow lane that the roadmap called for. One queued Staff 4 review packet from the local IMSLP Scherzo-Tarantelle solo part was converted into verified MusicXML for the visible E-flat5--E-flat5--C5--E-flat5--E-flat5 source measure-window. The accepted crop is a non-generated source-score image from page 2 with boxes centered on the five noteheads. The live May 3 detected phrase \texttt{D\#5 D\#5 C5 D\#5 D\#5} is accepted against that source window because the MIDI sequence \texttt{75 75 72 75 75} matches the verified score exactly; the displayed score spelling remains \texttt{Eb5 Eb5 C5 Eb5 Eb5}. The truth manifest now verifies two source-positive exact-MIDI phrases, marks one as live accepted, and keeps the three known false phrase maps blocked. The regression subset increased this lane to 34 focused tests covering the MusicXML slice, source-snippet gate, exact-MIDI acceptance, written-flat display, source-range verification, and rejected phrase regressions. This is still not long-phrase transcription: it is one verified source/audio phrase that can be extended outward into adjacent measures.

The 2026-05-17 anchored-expansion harness implements that outward-extension step as a gate rather than a manual judgment. The solve tracker now starts from the accepted Staff 4 source/audio anchor and tests adjacent left and right source-note ranges against available audio-agreed note runs. At the right-1 promotion stage, the six-note lane \texttt{Eb5 Eb5 C5 Eb5 Eb5 Eb5} was accepted, and the seven-note right-side lane \texttt{Eb5 Eb5 C5 Eb5 Eb5 Eb5 C5} became the next audit target. This prevents Curtis from treating a new random phrase candidate as progress while the current accepted lane has not yet extended cleanly.

The next 2026-05-17 expansion-search pass completes the first search-pool improvement requested by that gate. The expansion harness now searches raw detected note series from the daily record in addition to already-ranked score-candidate cards, records whether the best run came from the raw series or ranked group pool, reports raw-run counts in the tracker, and has a regression test proving that a raw exact seven-note Staff 4 run becomes \texttt{ready\_for\_truth\_review} rather than being missed. This still does not accept a longer phrase automatically; the expanded phrase remains hidden until exact MIDI, paired audio, source crop, and truth evidence all pass.

The 2026-05-17 audit-rejection pass turns the audited right-2 failure into a permanent negative example instead of leaving it as the next task. The manifest blocks four false phrase maps, including the Staff 4 source/audio expansion whose sixth note is \texttt{D5} against source \texttt{Eb5}. The expansion harness keeps that rejected case visible in the audit items but removes it from the current-best path. A later right-1 exact audio candidate for \texttt{Eb5 Eb5 C5 Eb5 Eb5 Eb5} now passes and is accepted; the useful constraint remains that the old seven-note expansion cannot be recycled as a solved phrase.

The 2026-05-17 adjacent-mining pass implements the next search step directly. Rather than waiting for a score candidate to happen to surface, Curtis now asks a narrow question over stored May 3 detected audio: does any window contain \texttt{75 75 72 75 75 75} or \texttt{75 75 72 75 75 75 72}? If yes, the window is exposed as an exact mining candidate for audit. If no, the tracker records the nearest stored miss and tells the next sync/transcription pass to mine more neighboring active windows. This keeps the accepted Staff 4 lane fixed and prevents the system from treating unrelated score matches as progress.

The 2026-05-17 03:23 AM EDT source-audio rescanner removes one bottleneck in that search. Curtis no longer depends only on note windows that happened to be stored before the current extractor and score-lane gates existed. It can now reopen the actual source media for the accepted Staff 4 anchor, extract a padded local audio window, run the current transcription pipeline over that audio, and merge the resulting note runs into the exact-MIDI expansion and mining harnesses. A rescanned source run is search evidence, not accepted evidence: if it disagrees with the verified source lane, it is reported as a blocker; if it agrees, it becomes an audit candidate, not a visible match.

The 2026-05-17 anchor-persistence pass prevents a fragile live rollback. The accepted Staff 4 anchor is now a truth-manifest object, not only a side effect of the current daily \texttt{matchGroups} display. If the display layer has zero accepted groups for May 3, the backend still resolves the verified source/audio anchor, uses its stored sample and timing to run source-audio rescanning, and feeds the same anchor into expansion and adjacent mining. The duplicate guard prevents double-counting when both the visible match group and the truth-manifest anchor are present.

The 2026-05-17 12:56 PM EDT tiled-rescan pass is the next implementation step from that state. The source-audio rescanner now covers more of the same accepted source sample by extracting five overlapping windows, not one. This increases the chance of finding the adjacent \texttt{Eb5} continuation in the original audio while preserving the exact-MIDI gate and keeping every mismatched tile candidate-only.

The 2026-05-17 01:43 PM EDT anchor-guided reproduction pass fixes the current-detector rollback exposed by the tiled scan. The broad transcription pass can still miss or segment the accepted Staff 4 anchor badly, but the rescanner now has a separate seeded check: use the accepted note-window timings, rerun current pYIN/YIN/spectral-onset agreement on those exact audio slices, and mark the anchor reproduced only when the exact MIDI sequence \texttt{75 75 72 75 75} is confirmed. Live verification reported five guided anchor events and \texttt{reproduced} anchor status, while the adjacent exact-MIDI search still found zero right-extension candidates after 1671 windows. The notation renderer was also tightened so accidental spelling affects display: \texttt{A\#} can render as \texttt{Bb} in a flat-key context or with a local accidental, rather than putting a natural \texttt{A} notehead on the staff. The 2026-05-17 02:19 PM EDT renderer patch moves local sharps and flats from text entities into drawn SVG accidental glyphs, keeps high-note accidentals inside the visible staff viewport, and verifies that \texttt{A\#}, \texttt{D\#}, \texttt{Bb}, \texttt{Eb}, and the \texttt{Bb}/\texttt{Eb} key signature all render with matching visible symbols.

The 2026-05-17 04:49 PM EDT adjacent-guided phrase pass is the next concrete implementation step from that state. Instead of only widening source windows again, Curtis now derives the Staff 4 right-1 and right-2 MIDI targets from the accepted source lane, infers the next note windows from the accepted anchor timing, and checks each expected note with current pYIN/YIN/spectral exact-MIDI votes. If every note reproduces, the longest adjacent target enters the hidden exact-MIDI mining pool as \texttt{staff4\_adjacent\_guided\_current\_detector}; if any note fails, the rescan record stores the first failed note, expected MIDI, detector votes, and seed window. The regression suite now proves that a broad transcription miss can still produce a seven-note adjacent phrase when current audio agrees with every source MIDI value, without allowing that phrase to bypass source crop and truth gates.

The 2026-05-17 05:29 PM EDT adjacent-onset sweep pass follows the same fixed-lane rule but removes a brittle timing assumption. Inferred continuation notes now sweep bounded early and late offsets around the predicted onset before failing. The accepted five anchor notes remain fixed; only notes beyond the anchor are swept. Each attempt records the offset, note window, exact detector sources, full detector votes, and best failed attempt. A new regression proves that when the base inferred sixth note returns \texttt{D5} but a nearby offset returns source \texttt{Eb5}, the seven-note Staff 4 target can enter the hidden exact-MIDI mining pool with the nonzero timing offset recorded.

The 2026-05-17 05:54 PM EDT failure-localization and notation-quality pass completes the next diagnostic layer. The source-audio rescanner now carries a compact first-failure object through the completion API and site tracker: direction, target sequence, failed note index, expected note and MIDI, best swept offset, detector consensus MIDI, and failure kind. This prevents the next pass from asking the broad question ``why no phrase?'' and instead asks the narrow question ``why did this exact Staff 4 adjacent note fail?'' The notation renderer also received a global symbol pass. Sharps, flats, key-signature flats, the treble clef, high accidentals, note labels, and CSS stroke/fill behavior are now covered by regression tests so accepted snippets cannot silently regress into crude or musically misleading glyphs.

The 2026-05-17 06:14 PM EDT notation visual-QA pass fixed the remaining symbol-rendering failure caught in browser review. The clef itself was no longer a note-name error, but the staff viewport still clipped the bottom of the glyph and compact responsive panels could shrink the notation into a lower-quality mobile rendering. Curtis now uses a taller notation viewport for both accepted transcription staves and notation gates, keeps compact matched panels from squeezing the clef, and preserves the same sharp/flat stroke and fill geometry across responsive overrides. Browser screenshots at desktop and mobile widths verify that the treble clef is not clipped, local sharps remain visible as sharps, local flats remain visible as flats, and the two-flat key signature renders on B-flat/E-flat staff positions.

The 2026-05-17 08:40 PM EDT Staff 4 first-failure audit gate adds the review loop needed before the lane can move again. The first failed adjacent-note object now preserves source spelling separately from detector spelling, so a Staff 4 source \texttt{Eb5} is no longer surfaced as \texttt{D\#5} just because the detector uses sharp names internally. The audit packet now records a decision: exact accepted source/audio evidence may extend the lane, wrong-MIDI detector consensus becomes a rejected regression case, and missing or non-reproduced detector evidence stays queued for gold review or rescanning. The site shows the decision label and source/audio note pair, which keeps the current work centered on the real question: whether the next adjacent Staff 4 source note is audibly and score-truth verified.

The 2026-05-17 10:46 PM EDT Staff 4 right-1 promotion converted the exact audio candidate into accepted source/audio truth. The accepted phrase was source \texttt{Eb5 Eb5 C5 Eb5 Eb5 Eb5}, detected audio \texttt{D\#5 D\#5 C5 D\#5 D\#5 D\#5}, exact MIDI \texttt{75 75 72 75 75 75}, local \texttt{20.225--22.992} inside \texttt{Njh8\_zq9\_DM-8835}, with the new actual-source crop \texttt{wieniawski-scherzo-tarantelle-staff4-eb-eb-c-eb-eb-eb-verified.png}. The truth manifest then had two live accepted Staff 4 source/audio phrases and four verified positive source phrases. The right-2 seven-note lane became the next audit target; it could not inherit acceptance from the six-note prefix.

The 2026-05-17 11:34 PM EDT anchor-advancement pass fixed the next source of circular work. The Staff 4 harness now prefers the longest accepted source/audio phrase that begins with the verified Staff 4 MIDI prefix, so the accepted base advanced instead of falling back to the older five-note anchor. Expansion moves rightward from that accepted lane; the previous left-side or shorter-prefix windows are no longer treated as the current task. The source target at that stage was exactly the seven-note lane \texttt{Eb5 Eb5 C5 Eb5 Eb5 Eb5 C5}. If audio evidence for that full MIDI sequence was found, it could become \texttt{ready\_for\_truth\_review}; it still could not display as accepted until the exact source crop, local audio/video, transcription glyphs, and truth gate all agreed. The stale sixth-note failure wording was also retired because the sixth note was accepted; the unresolved boundary then was the seventh source note, \texttt{C5}.

The 2026-05-17 11:42 PM EDT anchor-reproduction correction adds the accepted per-note audio windows to the truth manifest for the five-note and six-note Staff 4 phrases. The backend now consumes those windows as seed timing when the current daily display does not carry the accepted note events, and the source-audio rescan cache was advanced to \texttt{staff4\_source\_audio\_rescan\_v7} so live verification cannot reuse stale v6 anchor-reproduction records. This is not a new accepted phrase; it is the timing support needed for the current detectors to recheck the already accepted six-note anchor before testing the seventh-note C5 extension.

The 2026-05-17 11:51 PM EDT live audit verification was the endpoint state before truth promotion. The deployed API reported 69.113 percent implementation completion, reproduced the accepted six-note Staff 4 anchor under the v7 rescan, and found an exact seven-note MIDI candidate: source \texttt{Eb5 Eb5 C5 Eb5 Eb5 Eb5 C5}, detected audio \texttt{D\#5 D\#5 C5 D\#5 D\#5 D\#5 C5}, MIDI \texttt{75 75 72 75 75 75 72}. The Staff 4 audit packet \texttt{staff4-2026-05-03-Njh8\_zq9\_DM-8835-9-16} included playable local audio and video artifacts plus the source crop, but its status remained \texttt{pending\_source\_lock} and truth decision remained \texttt{pending\_review}. This was the first seven-note exact-MIDI Staff 4 candidate; it was not accepted evidence until the source crop/range and musician truth review were locked.

The 2026-05-18 Staff 4 full-phrase audit correction fixes a false bottleneck in that endpoint. The audit packet was carrying the exact seven-note stored audio run but using stale first-failure metadata, which made the generated playable review window cover only the first note. Version \texttt{staff4\_phrase\_audit\_v3} treats an exact current target/audio MIDI sequence as a full-phrase packet, keeps \texttt{firstFailure} empty, exposes a full-phrase check, and sets \texttt{pending\_source\_lock} rather than pretending the phrase is accepted. This means the next review can inspect the whole local \texttt{20.225--23.280} phrase window before any truth-manifest promotion.

The 2026-05-18 Staff 4 seven-note truth promotion completes that review step. The full phrase \texttt{Eb5 Eb5 C5 Eb5 Eb5 Eb5 C5} is accepted as source/audio truth against the detected audio spelling \texttt{D\#5 D\#5 C5 D\#5 D\#5 D\#5 C5}; both reduce to exact MIDI \texttt{75 75 72 75 75 75 72}. The accepted local window is \texttt{20.225--23.280} inside \texttt{Njh8\_zq9\_DM-8835}, and the accepted score crop is the same actual IMSLP Staff 4 source lane. The truth manifest now carries three live accepted Staff 4 source/audio phrases. The next implementation step was to verify one more source note and ask whether the original audio can support that extension.

The 2026-05-18 owner-media and A4-audit pass performs that next step without weakening the gate. Curtis now resolves \texttt{Njh8\_zq9\_DM-8835} from the local owner-media browser export when the state row is absent, seeds the accepted seven-note phrase from truth-manifest note windows, and probes only the new adjacent source note instead of retesting the accepted prefix. The v10 Staff 4 source-audio rescan produced six runs, 300 source-audio events, reproduced anchor status, two adjacent targets, and 258 searched adjacent windows. The first failed adjacent note is source \texttt{A4}; the best swept window is \texttt{23.514--23.677}, and current pYIN, YIN, and spectral-onset votes all report \texttt{D5}. The packet \texttt{staff4-2026-05-03-Njh8\_zq9\_DM-8835-9-17} now contains paired local audio/video artifacts and a \texttt{rejected\_mismatch} decision. The lane therefore remains seven accepted notes, and the next work is detector/segmentation calibration around the A4 continuation rather than another random score match.

The 2026-05-18 03:48 PM EDT failed-note pitch diagnostic turns that calibration target into measurable evidence. The Staff 4 audit packet now stores harmonic pitch-energy ranks for the exact failed-note window. On the current A4 continuation attempt, \texttt{D5} is the dominant candidate and \texttt{A4} carries only 0.276 of the observed \texttt{D5} harmonic energy, so the diagnostic class is \texttt{observed\_region\_dominates} and the packet stays rejected. A paired regression test proves the opposite case too: when the audio actually contains A4, the diagnostic reports \texttt{expected\_pitch\_present}. This is not a new phrase acceptance; it is the guard that separates true wrong-audio windows from future detector misses before Curtis tries to grow the lane again.

The 2026-05-18 04:10 PM EDT continuation-search pass is the next concrete implementation step from that guard. Curtis no longer tests only the timing inferred from the accepted Staff 4 anchor when the next source note fails. It now slides a short harmonic expected-pitch prefilter over the continuation region, keeps only the strongest candidate windows, and then applies the same exact pYIN/YIN/spectral-onset gate before adding any adjacent-guided audio run. The local v11 rescan found a hidden exact-MIDI eight-note Staff 4 candidate \texttt{Eb5 Eb5 C5 Eb5 Eb5 Eb5 C5 A4} / \texttt{75 75 72 75 75 75 72 69}; the added \texttt{A4} came from a later continuation window at local \texttt{32.265--32.425}, not from the rejected D5-dominant inferred window. This is now an audit target, not accepted display evidence: the next gate is to lock the expanded source crop, local media playback, rendered transcription, and truth review for the full eight-note phrase.

The 2026-05-18 04:43 PM EDT phrase-continuity gate rejects that hidden eight-note candidate as a stitched sequence rather than a phrase. Exact MIDI alone is no longer enough for phrase growth: the detected note windows must also be temporally continuous. The current Staff 4 eight-note run matches source MIDI, but its first seven notes occupy local \texttt{20.225--23.280} and the added \texttt{A4} is at \texttt{32.265--32.425}, leaving an \texttt{8.985} s internal gap. Curtis now marks this state as \texttt{blocked\_discontinuous\_audio\_phrase} / \texttt{exact\_midi\_discontinuous}, records a rejected regression in the truth manifest, keeps the eight-note source crop source-only, and extends the audit clip window so the late note is inspectable instead of silently omitted. The lane therefore remains seven accepted notes; the next transcription step is to find a continuous audio window for the same source notes, not to promote the stitched A4.

The 2026-05-18 05:19 PM EDT visible-mismatch quarantine supersedes the accepted Staff 4 lane. Alan reviewed the live 2026-05-03 score--clip--transcription group and reported that the visible score notes and transcription do not match. Curtis now treats that as a source-evidence failure, not a cosmetic issue: the five-note Staff 4 range \texttt{Eb5 Eb5 C5 Eb5 Eb5} and its dependent six-, seven-, and eight-note extensions are blocked by explicit rejected source ranges. The live deployment now has 0 live accepted Staff 4 phrases, 6 blocked regression cases, 0 visible May 3 match groups, 0 exact score-aligned windows, and 50.354 percent exact implementation completion. The source-audio rescan advances to \texttt{staff4\_source\_audio\_rescan\_v12}; its raw continuity-probe path can search a late-note neighborhood without truth-anchor or guided-note stitching, but it cannot run as accepted evidence until a new visible source-score crop is verified. The next step is therefore not a longer Staff 4 expansion; it is to reverify one May 3 source-score crop against the transcription and paired audio, then restart phrase growth only from that accepted anchor.

The 2026-05-18 06:14 PM EDT source-crop reverification queue implements that next step. A user-rejected Staff 4 visible mismatch no longer leaves Curtis with only a prose instruction or a stale hidden snippet. The completion API now exposes a concrete \texttt{queued\_source\_crop\_reverification} target for the May 3 lane: source score spelling \texttt{Eb5 Eb5 C5 Eb5 Eb5}, detected audio spelling \texttt{D\#5 D\#5 C5 D\#5 D\#5}, exact MIDI \texttt{75 75 72 75 75}, local sample \texttt{Njh8\_zq9\_DM-8835}, source window \texttt{*8835-8925}, the Staff 4 five-note crop, and note windows \texttt{20.225--22.535}. Running the Staff 4 audit generator now creates a \texttt{staff4-source-crop-review} packet for range \texttt{9--14} with playable local audio/video artifacts, the source crop under review, a pending review decision, \texttt{truthDecision=pending\_review}, and \texttt{canExtendStaff4Lane=false}. The 06:32 PM EDT packaged fallback ships the short review audio/video with Curtis so the live packet remains playable even if the live runtime volume cannot recreate it from owner media. This is not acceptance; it is the inspectable packet that must be approved before any Staff 4 anchor can return.

The 2026-05-18 07:56 PM EDT source/audio acceptance pass completes that small review step. The Staff 4 packet \texttt{staff4-source-crop-review-2026-05-03-Njh8\_zq9\_DM-8835-9-14} now accepts only the full-context actual-PDF crop for source \texttt{Eb5 Eb5 C5 Eb5 Eb5}. The old tight crop stays rejected. The score review verifies the visible source sequence and notehead centers; the audio review rechecks each stored note window with pYIN, YIN, and spectral peak and requires at least two agreeing detector sources per note. The accepted score MIDI and detected-audio MIDI both reduce to \texttt{75 75 72 75 75}. This restores one five-note source/audio anchor for Staff 4 expansion, but it does not restore the old six-, seven-, or eight-note extensions. Those dependent ranges must be audited again from the reverified anchor before Curtis may show them as accepted score--audio--transcription evidence.

The 2026-05-19 measure-16 display revocation replaces the earlier bracket-confirmation claim. Alan reviewed the live page and reported that the audio is measure 16, but the transcription should begin \texttt{Eb D C} rather than \texttt{Eb Eb C}, and the displayed crop is not the correct measure-16 score snippet. Curtis now records this as a failed acceptance gate. The five-note Staff 4 display, the red-bracket crop, and the measure-label filename fallback are blocked from accepted match display. Internal staff numbering is only extraction scaffolding; it is not a musical location and cannot substitute for source-score verification.

The 2026-05-19 Staff 4 six-note re-audit implements the next strict extension check without promoting a false phrase. Curtis now has a new full-context actual-PDF source review crop for \texttt{Eb5 Eb5 C5 Eb5 Eb5 Eb5}; the visible score notes and boxed notehead centers are treated as visually reverified source material. The paired audio gate still fails, so the phrase remains hidden from accepted match display. The packet \texttt{staff4-source-crop-review-2026-05-03-Njh8\_zq9\_DM-8835-9-15} stores playable audio/video, the source crop, pitch trace, spectrogram, and a blocked decision: notes 1--5 agree, but note 6 requires source \texttt{Eb5}/detected \texttt{D\#5} MIDI 75, while pYIN and YIN round that window to \texttt{D5} and only spectral peak returns \texttt{D\#5}. The current next step is therefore not a longer visible phrase; it is to find or segment a continuous six-note audio window where at least two independent detectors agree on the sixth \texttt{Eb5}/\texttt{D\#5}.

The 2026-05-19 failed-note raw-audio probe implements that next step as a reusable gate instead of another manual retry. When the adjacent-guided Staff 4 extension fails at the next source note, Curtis now opens a raw rescan centered on the failed local audio window, disables truth-anchor and guided-note seeding for that probe, and searches the resulting detected notes for the exact target MIDI sequence. If the probe finds a continuous exact-MIDI candidate, it can enter audit; if it does not, the current inferred adjacent window is locked as rejected and the next search must move to a different continuous attempt from the same measure-16 source lane. The source-crop six-note gate now runs the same failed-note raw-audio probe before repeating the current \texttt{Eb5 Eb5 C5 Eb5 Eb5 Eb5} window; the local completion API reports one probe, zero continuous candidates, and the next action to reject that window and search a different continuous measure-16 attempt. The follow-on 2026-05-19 alternate-attempt search turns that instruction into code: after the failed-note probe rejects the current timing, Curtis scans separate raw windows from the same May 3 source sample, excludes the known failed local start, and reports whether any alternate continuous exact-MIDI \texttt{75 75 72 75 75 75} candidate exists before another source-crop audit can be accepted. The current implementation now has a fourth larger-span gate: after the current 90-second source chunk fails, Curtis searches adjacent local owner-media chunks from the same May 3 video while keeping the same exact-MIDI and continuity requirements. Local verification now reports 12 alternate raw-window records, zero alternate continuous candidates, 22 wider overlapping source-window records, zero wide continuous candidates, and 28 adjacent-source records across \texttt{Njh8\_zq9\_DM-8745} and \texttt{Njh8\_zq9\_DM-8925}, with zero adjacent continuous candidates. The next action is therefore no longer to reuse the rejected window; it is to improve Staff 4 note segmentation or build a full-day Staff 4 attempt index. The source-audio cache remains \texttt{staff4\_source\_audio\_rescan\_v13}, and regression tests cover five outcomes: a raw probe that finds \texttt{75 75 72 75 75 75}, an alternate raw-window search that finds a continuous exact-MIDI candidate away from the failed start, a wider overlapping search that can find an exact-MIDI candidate missed by the sparse pass, an adjacent-source search that can find the phrase in a neighboring chunk, and a completion state that refuses to keep calibrating the same failed window when no alternate, wide, or adjacent continuous phrase is found.

The 2026-05-19 label-only verification failure is now a regression case. The previous deploy verified that the website label said \texttt{m. 16}, but that did not verify the music. The corrected gate now treats label correctness as insufficient: a displayed match needs the source crop, score notes, rendered transcription, paired audio, and measure location to agree. The front-end filename fallback for the Staff 4 measure-16 crop has been removed, and the rejected crop image is blocked at the source-image gate.

The 2026-05-20 02:15 AM EDT review-queue suppression update removes a training-loop burden exposed by live review. Curtis had been treating rejected review candidates as noisy soft labels, so exact normalized MIDI patterns and the same visible clip or score claim could reappear in the active queue and force Alan to reject the same wrong visual/audio claim again. The 2026-05-20 overlap-suppression follow-up closes the remaining loop: shifted windows from the same rejected source area are now hidden too. Live verification of the first overlap deployment then exposed the same burden on accepted areas: candidates marked \texttt{accepted\_candidate\_fingerprint} could keep reappearing from already accepted source windows. The gold-review learning profile now suppresses a rejected exact MIDI pattern, clip window, score-image fingerprint, overlapping rejected clip interval, or already accepted source area after one label. The front end also stops using recent reviewed items as fallback queue entries when the active queue is empty. Regression tests now prove that one rejected pattern is hidden from the queue, a same clip window with a different bad detected-note guess is hidden, overlapping shifted windows from a rejected area are hidden, accepted overlapping windows are treated as already covered, non-overlapping windows from the same source stay reviewable, and repeated rejected patterns stay suppressed while adaptive review mines fresh windows when the primary queue has been exhausted by rejections.

The 2026-05-20 original-score source-strip update addresses a separate evidence-presentation failure. The live site could show score-transcription fixtures and score-transcription queues, but it did not plainly show the rough original score material Alan expected. Curtis now crops six staff-level source snippets from page 2 of the local public-domain IMSLP Wieniawski Scherzo-Tarantelle solo-violin PDF and exposes them in the gold-review surface as original score source material. Those snippets are actual scanned score images, not generated notation. They remain source material only: they do not create an accepted score--audio match, practice evidence, or long-phrase transcription unless a later gate proves the specific source range, transcription, local audio/video, exact MIDI sequence, visible noteheads, and human review agree.

The 2026-05-20 requested-source library update closes the next presentation gap. Alan asked why the original score snippets in score-transcription did not look like original score material and why the requested repertoire list was missing. Curtis now keeps raw source metadata for the 13 requested repertoire targets: Haydn Symphony No. 94 IV violin I; Schubert Symphonies Nos. 3, 4, and 5 requested violin-I movements; Paganini Violin Concertos Nos. 1 and 2 solo movements; Paganini \emph{Moto perpetuo}; Wieniawski Etude-Caprices Op. 18 Nos. 3 and 4; Mozart \emph{Le nozze di Figaro} overture violin I; and Schumann Symphony No. 1 IV violin I. The visible score-transcription source list is stricter: it exposes only the ten entries backed by a rendered single-violin-part or solo-violin crop, and it hides complete scores, piano/accompaniment systems, orchestra sources, and multi-player source books from the review surface. Generated score-transcription packets are labeled as copy targets, not original score evidence. The visible source library is placed below the training queues as a compact score-PDF list rather than another review-card grid. This does not solve long-phrase transcription, but it removes a misleading evidence boundary: a source image shown as original score must come from a real violin source asset, and generated score-transcription notation must not masquerade as the piece.

The follow-up score-transcription wiring update fixed one review-card shape problem and exposed a deeper evidence problem. The earlier list exposed the source PDFs, but the score-transcription cards still rendered generated notation on both the left and right, which made the task look like Curtis was copying its own transcription. The score-transcription candidates then inherited source-library image paths, marked those images as original score snippets, and suppressed the generated source-notation fallback whenever a real source image was present. That solved the real-image requirement but still left false source/copy pairings: the displayed IMSLP crop and generated copy phrase were not verified as the same notes. Curtis now quarantines those generated requested-repertoire score-transcription packets until a verified source note sequence exists for the exact displayed crop. The source library still keeps the larger PDF/page preview and tighter review crops, but the review queue no longer asks Alan to accept or reject score-transcription pairs that Curtis has not proven are paired. The gold-review loop still tops up thin audio-review batches with adaptive audio windows, so the transcription-alan lane does not collapse while score-transcription extraction remains unresolved.

The same 2026-05-20 score-transcription gate now applies to older practice-day source snippets as well as the requested-repertoire source library. A source snippet no longer enters the score-transcription review lane merely because it has a real image, a measure label, an exact-note label, or an older accepted-truth flag. It must carry an explicit score-transcription-ready flag, a verified source-note sequence from the same visible crop, matching MIDI between that verified source sequence and the displayed copy target, no rejected/blocked/mismatch status, and no source-image-required fallback. This removes the remaining six stale Wieniawski score-transcription cards from the live review queue and prevents user-rejected Staff 4 and opening-crop variants from reappearing as copy tasks.

The 2026-05-21 verified score-transcription training update replaces that best-effort queue because Alan's review showed that plausible-looking source-copy cards were still mostly rejection work. Curtis now keeps the requested IMSLP source library visible as source material, but the \texttt{score-transcription} approval lane only exposes curated source-note packets whose copied pitch sequence is explicitly verified against the same visible source crop. The initial live packet set contains four actual IMSLP violin-part crops: a Mozart \emph{Figaro} five-note opening figure, the same Mozart opening D, the first two visible source notes from Schubert Symphony No. 4 movement IV rehearsal D, and the first visible source note from Haydn Symphony No. 94 movement IV measure 169. These packets are marked \texttt{verified\_source\_pitch\_copy\_training}, \texttt{sourceCopyReviewReady}, and \texttt{sourceCopyPitchSkeletonOnly}; they collect accept/reject labels for pitch, octave, accidental, order, and source-range agreement, not for final engraving-style equivalence. Unverified requested-repertoire crops no longer become score-transcription cards just because a real image exists. This deliberately reduces queue volume while raising expected approval rate: Curtis should show a few source-derived cards that Alan can likely approve, not many guessed cards that force repeated rejection. The accepted-evidence gate remains unchanged; none of these training packets can become accepted playing evidence, score-linked phrase evidence, or solved transcription without later audio, source-crop, exact-MIDI, continuity, and truth-review agreement.

The same 2026-05-21 review update adds a third training lane, \texttt{note-reading}, because Alan identified a simpler useful label: show a notation or original-score snippet and ask for note letters only, ignoring accidentals and octave numbers. The first \texttt{note-reading} queue is derived from the same four verified source crops as \texttt{score-transcription}, but it stores typed letter-sequence labels under \texttt{note\_letter\_reading}. A same-day correction changed the lane from a quiz to a human-labeling tool: a submitted sequence such as \texttt{d c d c d} is Alan's accepted source-note label, and disagreement with Curtis' current expected sequence is stored as correction signal rather than displayed as \texttt{Failed}. These labels do not create accepted playing evidence, accepted score evidence, or solved transcription. They isolate one subproblem that Curtis kept forcing into harder accept/reject cards: whether the system and reviewer agree on the basic note-letter skeleton visible in the source notation. This should make review less brittle and produce higher-volume training data before Curtis asks Alan to judge octave, accidental spelling, rhythm, score range, and audio agreement together.

The 2026-05-21 review-playback cleanup removes another labeling burden from the same training surface. The \texttt{note-reading} input no longer shows a fake example such as \texttt{a g b c d}; an empty field now means Alan is supplying the note-letter label for the displayed source picture only, not hidden surrounding notes. \texttt{transcription-alan} and score-linked review packets also no longer cut playback exactly at the detector note boundary. Curtis stores the exact note-local start and end times separately, but the playable audio/video review window now carries 0.4 s of pre-roll and 0.65 s of post-roll when media limits allow. This makes accept/reject labels less vulnerable to clipped attacks, clipped releases, and source-score packets whose audible note begins just before the old review window.

The remaining M1 work is continuing the low-cost media pass over every unmeasured strict-ledger video in chronological windows, then replacing the sample-biased estimate with full-archive checked intervals. The remaining M10 work is the larger benchmark suite for repeated notes, arpeggios, fast runs, score-free exercises, score-coordinate truth, accepted notation/media completeness, and practice-time accounting regressions.

The long-phrase track should now run in parallel with full-archive scanning. Waiting for all 276 h 23 min of archive coverage before extending one musical lane would be too slow. The faster strategy is now a narrow reverified-measure lane:
\begin{enumerate}
    \item pick one May 3 local audio/video phrase and its proposed Scherzo-Tarantelle source range;
    \item verify the visible MusicXML notes, actual-score crop, boxed notehead centers, transcription, and paired audio before any display acceptance;
    \item reject or accept that source/audio crop as the new anchor;
    \item only after acceptance, search adjacent stored and source-rescanned detected note series against that same source lane;
    \item accept the expansion only if the exact score notes, displayed transcription, local audio, local video, source crop, and truth gate all agree;
    \item extend outward one verified adjacent note group at a time until the lane becomes a phrase.
\end{enumerate}

The next coding pass should:
\begin{enumerate}
    \item create a new verified May 3 source--audio--transcription pair before any current anchor can return;
    \item keep the old tight crop and all dependent six-/seven-/eight-note ranges blocked until each receives a fresh exact source/audio audit;
    \item keep the raw continuity-probe path available for future accepted anchors, but do not use it to revive rejected extensions automatically;
    \item search stored and source-rescanned audio for a continuous adjacent phrase only after a new accepted anchor exists;
    \item accept a longer compact grouped score--clip--audio--transcription pair only if the expanded exact MIDI sequence, paired audio, source crop, temporal-continuity gate, and truth gate all pass;
    \item keep reporting the first blocking note when an expansion fails so the next extraction attempt has a precise target;
    \item keep running queue-only owner-media sync in parallel, starting with \texttt{violin 1} at \texttt{*2250-2340};
    \item rerun the active-practice scan after those uploads and record checked, active, and pending totals.
\end{enumerate}

Curtis currently has no accepted Staff 4 visible match after the live measure-16 review. It cannot claim a longer score match from guessed scanned crops, one-note anchors, reference-audio pitch traces, source notes that have not been verified into MusicXML, user-rejected visible crops, label-only measure matches, internal staff indices, or exact-MIDI sequences stitched across separate playing moments. The speed-up is to find one new exact score--audio--transcription pair, then grow only from that accepted pair.
